\documentclass[10pt]{article}
% \documentclass[10pt, draft]{article} % draft mode (no figures)

\usepackage[margin=1in]{geometry}
\usepackage{amsmath,amssymb,amsthm,bm,mathtools}
\usepackage{algorithm}
\usepackage{algpseudocode}
\usepackage{enumitem}
\usepackage{graphicx}
\usepackage{subcaption}
\usepackage{booktabs}
\usepackage{tabularx}
\usepackage{makecell}
\usepackage{float}
\usepackage{xcolor}
\usepackage{hyperref}
\usepackage[style=numeric,sorting=nyt]{biblatex}
\usepackage{dsfont}

\usepackage{color,soul}
\usepackage[textsize=tiny]{todonotes}
\usepackage[normalem]{ulem} % provides \sout

\definecolor{impGreen}{RGB}{0,128,0}



\addbibresource{references.bib}
\addbibresource{references_additions.bib}

\hypersetup{hidelinks}

\newtheorem{theorem}{Theorem}
\newtheorem{proposition}{Proposition}
\newtheorem{lemma}{Lemma}
\newtheorem{corollary}{Corollary}
\theoremstyle{definition}
\newtheorem{definition}{Definition}
\newtheorem{assumption}{Assumption}
\newtheorem{remark}{Remark}

\newcommand{\R}{\mathbb{R}}
\newcommand{\E}{\mathbb{E}}
\newcommand{\Cov}{\operatorname{Cov}}
\newcommand{\Corr}{\operatorname{Corr}}
\newcommand{\Var}{\operatorname{Var}}
\newcommand{\diag}{\operatorname{diag}}
\newcommand{\1}{\mathbf{1}}
\newcommand{\D}{\mathcal{D}}
\newcommand{\Hcal}{\mathcal{H}}
\newcommand{\Vcal}{\mathcal{V}}
\newcommand{\argminop}{\operatorname*{arg\,min}}
\newcommand{\argmaxop}{\operatorname*{arg\,max}}
\newcommand{\advisorchange}[1]{{\color{blue}#1}}
\newcommand{\blueheading}[1]{\texorpdfstring{\textcolor{blue}{#1}}{#1}}
\newenvironment{advisorchanges}{\begingroup\color{blue}}{\endgroup}

% Allows the source to compile outside the project directory.
\newcommand{\safeincludegraphics}[2][]{%
  \IfFileExists{#2}{\includegraphics[#1]{#2}}{%
  \fbox{\parbox[c][0.19\textheight][c]{0.92\linewidth}{\centering
  Figure file not found in this folder.\\[2mm]{\scriptsize\ttfamily\detokenize{#2}}}}}}

\title{Reducing Regressivity in ML-Based Mass Appraisal:\\
A Covariance-Guided Correction for Assessor Workflows}

\author{
Nicol\'as Acevedo\thanks{Operations Research Center, MIT, Cambridge, MA 02139. \texttt{nacevedo@mit.edu}}
\qquad
Haihao Lu\thanks{Sloan School of Management, MIT, Cambridge, MA 02139. \texttt{haihao@mit.edu}}
\\[1ex]
Michael Wagner\thanks{Cook County Assessor's Office, Chicago, IL 60602. \texttt{Michael.Wagner2@cookcountyil.gov}}
\qquad
Timothy Sparer\thanks{Cook County Assessor's Office, Chicago, IL 60602. \texttt{Timothy.Sparer@cookcountyil.gov}}
}
\date{}

\begin{document}
\maketitle

\begin{abstract}
\begin{advisorchanges}
Property value assessments determine how local property-tax burdens are distributed. Assessor offices increasingly use machine learning-based automated valuation models to estimate values for large and varied property inventories, but strong average prediction does not prevent \emph{regressivity}: lower-value properties can receive systematically higher assessment-to-sale ratios than higher-value properties. Working with data scientists and assessors at the Cook County Assessor's Office (CCAO), we develop a correction that can be added to its existing LightGBM residential workflow without replacing its data, features, or review process. The method adds a training penalty for a systematic relation between assessment errors and sale prices. We consider a direct squared-covariance penalty and a sample-additive, covariance-guided objective that is compatible with standard boosting software. Applied to a Python translation of the CCAO workflow, the validation-selected sample-additive model substantially improves three standard measures of vertical equity: their departures from neutral values fall by about 54\% to 88\% relative to ordinary LightGBM. Predictive fit changes little, while mean absolute error rises by about 3.5\%. The current comparison measures the full implemented custom-objective path and does not yet isolate the penalty from every implementation difference. The correction targets a global first-order pattern. It does not set the overall assessment level, remove local inequity, or guarantee that every ratio-study measure will fall within its reference range. It should therefore be used with temporal validation, uncertainty analysis, subgroup ratio studies, and comparison with simpler post-hoc recalibration.
\end{advisorchanges}
\end{abstract}

\section{Introduction}
\label{sec:introduction}

Property taxes are one of the largest sources of local government revenue in the United States. They help finance public schools, roads and streets, and other essential local services \cite{IAAO2020PropertyTaxPolicy}. Using 2016 U.S. Census Bureau data, \textcite{Berry2021} reports that property taxes accounted for roughly 72\% of local tax revenue and 47\% of local own-source general revenue, and totaled about \$500 billion annually.

% \textcolor{blue}{\st{The U.S. property tax system is generally value-based}}\textcolor{impGreen}
{Property taxation in the United States is generally based on property values}: local assessors estimate property values, which are then used, together with jurisdiction-specific assessment and taxation rules, to determine each taxpayer's share of the property-tax burden \cite{IAAO2020PropertyTaxPolicy,IAAO2025Mass}. Consequently, systematic overassessment of some properties and underassessment of others can shift tax liability from one group of property owners to another, creating or exacerbating inequities \cite{Berry2021,IAAO2020PropertyTaxPolicy}.

The {valuation} of groups of properties for tax purposes is known as \emph{mass appraisal}. The International Association of Assessing Officers (IAAO) defines mass appraisal as ``the process of valuing a group of properties as of a given date using common data, standardized methods, and statistical testing'' \cite{IAAO2025Mass}. Within this process, assessors use statistical models to estimate property values across large property groups. {When implemented in computer-based valuation systems, such models are commonly referred to as automated valuation models (AVMs)} \cite{IAAOAITaskForce2022,BidansetRakow2022}. Recent mass-appraisal research increasingly examines nonlinear machine-learning (ML) methods, including gradient-boosted trees, alongside traditional regression approaches \cite{Rootetal2023,Sevgen2024100111,Zilli2024}. %\textcolor{blue}{\st{Given their high performance in predictive power.}}

However, when evaluating the performance of a mass-appraisal valuation model, predictive accuracy alone is insufficient. The IAAO recommends the use of \emph{ratio studies}, which compare {appraised values} with observed sale prices. 
% \textcolor{blue}
% {\st{A central quantity in these studies is the assessment-to-sale ratio, defined as the assessed value divided by the sale price.}} 
{A central quantity in these studies is the ratio of appraised value to sale price. Because the CCAO model studied here predicts \emph{fair market value} before the statutory conversion to assessed value, we evaluate its predictions on a common market-value scale, using the model-estimated market value divided by the observed sale price; we refer to this quantity as the \emph{assessment ratio} throughout the paper \cite{CCAOModelResAVM2026}.} Ratio studies use such ratios to evaluate assessment level and uniformity, including whether similar properties are treated consistently (\emph{horizontal equity}) and whether assessment ratios vary systematically across property-value levels (\emph{vertical equity}) \cite{IAAO2013Ratio,IAAO2025Mass}.

{A particularly relevant form of vertical inequity is} assessment \emph{regressivity}: lower-value properties tend to receive higher assessment ratios than higher-value properties. This pattern can place a disproportionate share of the assessment, and potentially tax, burden on owners of lower-valued properties \cite{McMillen2020,Berry2021,AvenancioLeonHoward2022}. Moreover, regressivity may remain hidden in aggregate measures of predictive accuracy because these measures do not capture whether assessment errors vary systematically with property value. A valuation model can therefore achieve strong predictive accuracy while still exhibiting systematic price-related bias and, consequently, inequitable assessment outcomes.

% IS THIS PARAGRAPH WORTH KEEPING???
% \textcolor{blue}{\st{For property valuation modeling, this means that regressivity cannot be treated only as an ex-post diagnostic. If the goal is to preserve predictive performance while reducing systematic price-related bias, vertical equity must also enter model training or calibration as an explicit criterion. Implementing this objective is especially challenging in the mass-appraisal context: current sale prices are observed only when properties transact, so values for all properties in the jurisdiction must be inferred from a smaller and selected set of verified sales. Properties differ in size, age, condition, construction, use, and amenities; location effects operate at multiple geographic levels and vary over time; and administrative records may be missing, outdated, or inconsistent. At the same time, assessment systems must produce values within fixed assessment schedules while remaining auditable and defensible in taxpayer appeals. These constraints limit the types of equity corrections that are useful in practice: a method should reduce regressivity while preserving predictive performance, require limited changes to existing valuation workflows, and remain feasible to implement, audit, and defend within assessor timelines. Moreover, because valuation practices, available data, and administrative requirements differ across jurisdictions, the specific implementation of such a correction should also be adaptable to the local assessment context.}}
% \textcolor{blue}{\cite{IAAO2020SalesVerification,IAAO2018AVM,IAAO2025Mass,Almy2014}}

{This motivates treating regressivity as an explicit criterion in model training or calibration rather than only as an ex-post diagnostic. In mass appraisal, a useful correction must operate within existing sales-based modeling, validation, and review processes: current prices are observed only for properties that transact, while valuation systems must produce values within fixed assessment schedules and remain compatible with professional and administrative review \cite{IAAO2020SalesVerification,IAAO2018AVM,IAAO2025Mass}. Accordingly, the correction should reduce regressivity while preserving predictive performance and requiring minimal changes to existing valuation workflows.}

This paper arises from a collaboration between the Cook County Assessor's Office (CCAO) and academic researchers, motivated by an operational goal: reducing regressivity in the CCAO's residential valuation workflow. The CCAO currently uses a county-wide LightGBM (Light Gradient-Boosting Machine) model \cite{lightgbm} as its primary residential valuation model, {following earlier township-level linear-regression and gradient-boosting workflows} \cite{CCAOModelResAVM2026}. The current model is trained on screened residential sales and estimates a \emph{fair market value} using property characteristics, location, environmental variables, and market trends. It operates within a broader assessment pipeline that separates data ingestion, model training, assessment, evaluation, interpretation, and finalization \cite{CCAOModelResAVM2026}.

% \textcolor{blue}{\st{Further, the value produced by the valuation model is only one input into the final property-tax bill. In Cook County, the Assessor's estimate of fair market value is first converted to an assessed value according to property classification, for example residential or commercial. A state equalization factor is then applied, and eligible exemptions can further reduce the resulting equalized assessed value. Local tax rates are ultimately applied to this taxable equalized assessed value to determine the tax bill. These rates are calculated by the Cook County Clerk based on the tax levies submitted by the relevant taxing districts and the taxable property values within those districts. In this paper, we focus solely on the valuation stage of the property-tax system, rather than the subsequent steps that convert estimated market values into final tax bills.}}
% \textcolor{blue}{\cite{CCAOAssessmentNotice2026,CCAOPropertyTaxSystem2026}}

% \textcolor{impGreen}
{The model's fair-market-value estimate is only one input into the final property-tax bill. In Cook County, that value is converted to assessed value according to property classification, equalized using the state equalization factor, reduced by eligible exemptions, and ultimately combined with local tax rates determined from taxing-district levies and taxable property values \cite{CCAOAssessmentNotice2026,CCAOPropertyTaxSystem2026}. We therefore focus on the valuation stage rather than the downstream steps that convert estimated market values into final tax liabilities.}

Within the valuation stage, the central problem addressed in this paper is the interaction between predictive performance and {regressivity}. In the CCAO workflow considered, the gradient boosting model substantially improves predictive accuracy and ratio dispersion relative to {the linear-regression baseline}, but its vertical-equity diagnostics move in a more regressive direction. {This creates a model-design problem: rather than choose between the linear and gradient-boosting baselines, we seek to retain the predictive advantage of gradient boosting while explicitly reducing regressivity.} From an operational perspective, {the correction should also fit within CCAO's existing modeling and validation workflow without requiring} a fundamental change to CCAO's existing modeling and validation workflow \cite{CCAOModelResAVM2026}.

% \textcolor{blue}{\st{Within the valuation stage, the central problem addressed in this paper is the interaction between predictive performance and vertical equity. In the CCAO workflow considered, the gradient boosting model substantially improves predictive accuracy and ratio dispersion relative to linear regression-based models, but its vertical equity diagnostics move in a more regressive direction. This tradeoff between a simpler linear model and a more complex nonlinear model moves the discussion from a model-selection problem to a model-design problem: the goal is to preserve the predictive advantage of the nonlinear model while explicitly controlling its regressive behavior. From an operational perspective, we aim to achieve this goal without requiring a fundamental change to CCAO's existing modeling and validation workflow.}}
% \textcolor{blue}{\cite{CCAOModelResAVM2026}}

We address this challenge by modifying the training objective of the valuation model so that {regressivity} becomes an explicit model-design consideration. {Specifically, we formulate a direct squared-covariance penalty that controls global price-related dependence between valuation errors and property values, and then develop a sample-additive surrogate motivated by the same dependence-control objective that is compatible with LightGBM's custom-objective interface.} {For a fixed penalty formulation, the correction changes only the training objective and introduces one additional tuning parameter: the penalty strength.} The correction requires no changes to the underlying data, feature set, or surrounding validation and evaluation workflow. The {penalty strength} can be selected through chronological validation using the predictive and ratio-study measures already used to evaluate mass-appraisal models. The resulting approach is therefore designed to reduce regressivity while remaining compatible with the existing machine-learning valuation pipeline. {The framework is intentionally targeted: it does not, by itself, determine overall assessment level or guarantee local or subgroup equity.}

Applied to CCAO's residential property-valuation workflow, the proposed methodology substantially reduces measured regressivity while preserving nearly all of the predictive performance of the unpenalized LightGBM. We use the oldest 90\% of eligible 2016--2024 sales (344,610 sales) for model development, evaluate the selected models on the newest 10\% (38,290 sales) as a held-out test set, and use 2025 (26,641 sales) as {a secondary forward evaluation aligned with the assessment-year workflow} after refitting on all 382,900 eligible 2016--2024 sales. Relative to unpenalized LightGBM, the validation-selected model reduces the deviations of {three complementary vertical-equity diagnostics} from their neutral values by approximately 54\% to 88\%. At the same time, $R^2$ decreases only from 0.885 to 0.881, while mean absolute error increases by approximately 3.5\%. Thus, the proposed methodology produces a {substantial reduction in measured regressivity} at a small cost in predictive accuracy. Exploratory experiments across six counties further {suggest} that the correction can improve value-related assessment patterns in different markets, although the appropriate penalty strength and attainable improvement vary across jurisdictions. \todo{Update all numerical results after the final experimental run. Add confidence intervals? Is mean enough?}

The method was developed in collaboration with CCAO data scientists and has been incorporated into the office's residential model-development codebase. \cite{CCAOModelResAVM2026}\textcolor{impGreen}{\todo{Replace or supplement this citation with the specific CCAO commit/PR containing the correction.}} The CCAO has also used and tested the correction during model development. This implementation provides direct evidence that the proposed approach can be integrated into an existing public-sector mass-appraisal workflow without requiring a fundamental redesign of the underlying valuation pipeline. 

% \textcolor{blue}{\st{Accordingly, the method should be understood as targeting vertical equity in the valuation stage rather than equity in final property-tax liabilities.}}

{
\color{blue}
The \todo{OLD version: safely delete(?)} paper makes three contributions. First, motivated by an important operational problem recognized by CCAO, \textcolor{blue}{\st{we translate property-tax regressivity into a tractable model-training objective that explicitly targets the relationship between valuation errors and property values.}}}
\textcolor{impGreen}{we connect assessor-defined assessment regressivity with dependence-based constrained learning by translating the value-related error pattern into a tractable model-training objective.} Second, \textcolor{blue}{\st{we develop an easy-to-incorporate implementation that operates through the standard LightGBM interface and fits within an assessor's existing temporal-validation and ratio-study workflow.}}\textcolor{impGreen}{we develop a direct squared-covariance penalty and a sample-additive surrogate that can be implemented through the standard LightGBM interface and evaluated within an assessor's existing temporal-validation and ratio-study workflow.} Third, we demonstrate that the correction can substantially reduce regressivity while largely preserving predictive accuracy, and \textcolor{blue}{\st{we document its progression from a jointly developed research method to a component of CCAO's core residential valuation codebase.}}\textcolor{impGreen}{we document its incorporation into CCAO's residential model-development codebase.}
% \todo{I think we need to discuss this further. I see as a contribution: (1) a bridge between assessors literature and fairness in ML, (2) and implementation in a public workflow proved. }

% \color{red}
% The paper makes three main contributions. First, we connect assessor-defined {assessment regressivity} with dependence-based constrained learning by formulating assessment regressivity as systematic dependence between valuation errors and property value. This provides a direct way to translate a standard mass-appraisal equity concern into an explicit model-training objective. Second, {we develop a direct squared-covariance penalty and a sample-additive surrogate motivated by the same dependence-control objective, and show how these formulations can be incorporated into regression models used in mass appraisal, including linear regression and gradient boosting. The penalty strength is selected through chronological validation using predictive and ratio-study criteria.} Third, using CCAO's residential valuation workflow, we show that the proposed correction can substantially reduce measured regressivity with limited loss in predictive accuracy. The method has also been implemented and tested within CCAO's model-development pipeline, providing evidence of its operational compatibility with a large public-sector mass-appraisal system.
% \todo{I think we need to discuss this further. I see as a contribution: (1) a bridge between assessors literature and fairness in ML, (2) }
% }

% \textcolor{blue}{\st{It is important to highlight that the proposed correction is intentionally targeted. It controls a global first-order relationship between valuation errors and property value, but does not determine overall assessment level, eliminate local or subgroup inequities, or guarantee that all ratio-study diagnostics fall within their reference ranges. It is therefore intended to complement, rather than replace, improvements in data, local market modeling, calibration, and professional review.}}

% {\color{gray}
Our contribution is application-driven. Motivated by the regressivity problem observed in the CCAO residential valuation workflow, we study whether this problem can be addressed through the model-training objective rather than by redesigning the surrounding valuation architecture. We translate this regressivity problem into a first-order dependence between log valuation residuals and log sale prices, and penalize this dependence during model training using covariance-based regularization. To make this practical for nonlinear models used by assessors, we characterize the cross-observation coupling induced by the direct penalty and develop a sample-additive surrogate, motivated by a covariance upper-bound argument, that can be implemented in standard gradient-boosting interfaces. We evaluate these corrections \todo{Future note: this assumes that we are running the results in their pipeline...} within the CCAO residential valuation workflow and show, using chronological out-of-sample evaluation and standard ratio-study diagnostics, that they can substantially reduce regressivity while retaining similar predictive performance. Analytical results clarify the mechanism and resulting accuracy--equity tradeoff, while supplementary multi-jurisdiction experiments examine the extent to which the same objective-level correction transfers beyond Cook County and where its first-order formulation is insufficient.

% }

{The remainder of this section reviews the related literature on machine learning in mass appraisal, vertical equity, and dependence-based fair regression. The rest of the paper is organized as follows.} Section~\ref{sec:setting} describes the mass-appraisal setting, the CCAO valuation workflow, and the evaluation framework. Section~\ref{sec:method} develops the proposed covariance-based regressivity correction. The remaining sections present the empirical evaluation, examine the robustness and practical implementation of the approach, and discuss its limitations and extensions.

\todo{I am just missing how to use the figure in the regressivity part of the introduction}


% % \newpage
% \subsection{Related Work}
% \label{subsec:related_work}

% % \paragraph{Machine learning in mass appraisal.}
% % Research and assessor practice have moved from traditional regression toward tree ensembles and other nonlinear methods \cite{WangLi2019,Rootetal2023,Sevgen2024100111,Zilli2024}. Their use does not remove the need for sales verification, temporal testing, ratio studies, interpretation, and administrative review \cite{IAAO2020SalesVerification,IAAOAITaskForce2022,BidansetRakow2022}. The CCAO pipeline follows this sequence: model training is followed by assessment, evaluation, interpretation, and final review \cite{ccao_data_2026}.

% % {
% % \color{red}
% \paragraph{Machine Learning in Mass Appraisal} Recent reviews document the growing use of automated valuation models and machine learning (ML) in mass appraisal and real-estate valuation. Multiple regression-based models remain common, alongside with tree ensembles (e.g., Random Forest and Gradient Boosting), and neural-network-based methods \cite{WangLi2019,Rootetal2023,Zilli2024,purdueAssessingTechniques}. Most studies compare these methods primarily through predictive criteria, such as out-of-sample accuracy and performance relative to alternative algorithms. On the other hand, equity is less often incorporated directly into model training, model selection, and performance assessment \cite{WangLi2019,Rootetal2023,Zilli2024}. In terms of predictive  power alone, tree ensembles are among the most frequently studied and often strongest-performing nonlinear methods in this literature \cite{Rootetal2023,Sevgen2024100111,Zilli2024}. Therefore, a practical equity-aware intervention and assessment should be compatible with these already widely used model classes rather than restrict solutions to a specialized class of valuation models.

% Even though adoption of ML models is increasing, it is important to note that their adoption by assessment offices remains partial and uneven. Among the perceived adoption barriers, \textcite{BidansetRakow2022} identify
% factors such as: communication challenges, limited technical capacity, and budget constraints. Further, for assessor's offices, the predictive modeling stage is only one part, of many, involved in the assessment process; the use of ML does not replace sales verification, out-of-time evaluation, ratio studies, model interpretation, administrative review, and further adjustments \cite{IAAO2020SalesVerification,IAAOAITaskForce2022}. As one of the main assessors organizations, the \textcite{IAAOAITaskForce2022} emphasizes explainability, defensibility, administrative feasibility, and integration into existing workflows as central conditions for operational use of advanced ML models. These considerations favor equity-aware methods that can be incorporated into already existing regression objectives and model-development pipelines used by assessors, with limited disruption. This practical requirement motivates the approach developed in this paper.
% % }

% % \paragraph{Vertical equity and regressivity.}
% % \advisorchange{The assessor literature distinguishes horizontal dispersion from systematic value-related bias \cite{IAAO2013Ratio,cornia2006property,McMillenSingh2023}. Assessor standards use several complementary diagnostics because no single statistic fully describes vertical equity. Large samples can hide local pockets of bias, and regression-based diagnostics can behave differently when assessed values are themselves produced by regression models \cite{IAAO2013Ratio,IAAO2023VerticalReview,McMillenSingh2023}.} Proposed corrections include spatial and geographically weighted methods, stratified procedures, ratio-study-based validation, and fairness-aware valuation models \cite{Quintos2014,BidansetEtAl2019,BidansetEtAl2025,HermansEtAl2022,HermansEtAl2025,CandoganHanLu2023}. \advisorchange{The correction developed here adds a global price-related target during training} and leaves the standard ratio-study review in place.

% % {\color{red}
% \paragraph{Vertical equity and regressivity.}
% Equity has long been a central objective in mass appraisal practice and standards. The assessor literature distinguishes \emph{horizontal equity}, which concerns the uniform treatment of similar properties, from \emph{vertical equity}, which concerns systematic differences in assessment outcomes across property-value levels \cite{IAAO2013Ratio,cornia2006property,McMillenSingh2023}.  A common form of vertical inequity is \emph{regressivity}: assessment-to-sale ratios tend to decline with property value, so lower-valued properties tend to be relatively over-assessed and higher-valued properties tend to be relatively under-assessed. Ratio studies evaluate these patterns through stratified analyses and complementary measures of uniformity and price-related bias, such as the Coefficient of Dispersion (COD), the Price-Related Differential (PRD), and the Price-Related Bias (PRB) \cite{IAAO2013Ratio,IAAO2023VerticalReview}.

% From the ratio studies measures definitions, one can note that no single statistic fully characterizes vertical equity. Aggregate results of these metrics can conceal bias within geographic or property subgroups, and some regression-based diagnostics can behave poorly, in terms of equity quantification, when estimated values are generated by regression models \cite{IAAO2023VerticalReview,McMillenSingh2023}. Accordingly, assessor guidance favors examining several diagnostics all together, instead of relying on single global measures. Empirical studies further show that regressivity is widespread, that its consequences are of substantial importance, and that limitations in assessment data and valuation models themselves are among its likely causes \cite{Berry2021,AvenancioLeonHoward2022}.

% A smaller body of literature moves beyond measurement toward explicit correction of inequities within the valuation models. Proposed approaches include spatial and geographically weighted methods, stratified procedures, ratio-study-based model validation, and fairness-aware valuation models \cite{Quintos2014,BidansetEtAl2019,BidansetEtAl2025,HermansEtAl2022,HermansEtAl2025,CandoganHanLu2023}. The method developed here contributes to this strand by adding a global price-related equity-aware target during training procedures, while retaining standard ratio studies and subgroup review for its out-of-sample evaluation.
% % }

% % \paragraph{Fair regression and dependence penalties.}
% % Fair-regression methods use constrained estimation, post-processing, or dependence penalties for continuous outcomes \cite{AgarwalDudikWu2019,FitzsimonsEtAl2019,ChzhenEtAl2020Wasserstein,ChzhenEtAl2020Recalibration,WeiEtAl2023}. Prior work has used covariance, correlation, kernel dependence, and maximal-correlation penalties \cite{KomiyamaEtAl2018,MaryEtAl2019,PerezSuayEtAl2017,LiEtAl2022Fairness,ScutariPaneroProissl2022,lee2022maximal}. We use that machinery for an assessor-specific target: the value-related pattern in assessment ratios. The paper then studies implementation, temporal calibration, and failure modes in mass-appraisal models.

% % {
% % \color{red}
% \paragraph{Fair regression and dependence-based constrained learning.}
% Although much of the fairness-in-ML literature initially focused on classification, \emph{fair regression} studies how fairness requirements can be incorporated when predictions are continuous, as it is in mass appraisal. Existing approaches include constrained or regularized empirical risk minimization and post-processing methods. These methods impose fairness-aware criteria to the ML pipelines, such as statistical parity, bounded group loss, distributional parity through optimization hard constraints, soft-constraint through penalties, and recalibration procedures \cite{AgarwalDudikWu2019,FitzsimonsEtAl2019,ChzhenEtAl2020Wasserstein,ChzhenEtAl2020Recalibration,WeiEtAl2023}. The broader methodological principle is that predictive error loss (e.g., sum of the squares of the residuals) does not need to be the only single training objective: an application-specific requirement can be added as a hard constraint or soft penalty to the learning procedure, which in most cases produces an explicit tradeoff between predictive performance and the relevant notion of fairness.

% A particularly relevant literature for this paper represents unfairness as statistical dependence between model outputs and a variable of concern. Prior work has controlled this dependence using covariance or correlation, coefficients of determination, or kernel-based measures such as Hilbert-Schmidt Independence Criterio (HSIC), and maximal correlation \cite{KomiyamaEtAl2018,MaryEtAl2019,PerezSuayEtAl2017,LiEtAl2022Fairness,ScutariPaneroProissl2022,lee2022maximal}. These methods are especially useful when the relevant variable is continuous and can often be incorporated into standard regression objectives. However, our application differs from conventional group-fairness settings: property value is not treated as a protected attribute since it is the target itself. Instead, we use dependence regularization to encode an assessor-specific structural requirement by penalizing the property value-related pattern in assessment ratios associated with regressivity. In this sense, the proposed method is a form of domain-informed constrained learning that translates a standard ratio-study concern into a trainable objective.
% % }







\subsection{Related Work}
\label{subsec:related_work}

\paragraph{Machine Learning in Mass Appraisal}
Recent reviews document the {growing research on} automated valuation models and machine learning (ML) in mass appraisal and {real estate} valuation. {Multiple regression remains common, alongside tree ensembles such as random forests and gradient boosting, as well as neural-network-based methods} \cite{WangLi2019,Rootetal2023,Zilli2024,purdueAssessingTechniques}. Most studies compare these methods primarily through predictive criteria, such as out-of-sample accuracy and performance relative to alternative algorithms. {Vertical equity, although routinely evaluated in assessor practice, is much less often incorporated directly into model training or model selection} \cite{WangLi2019,Rootetal2023,Zilli2024}. {When predictive performance is the primary criterion, tree ensembles are among the most frequently studied and often strong-performing nonlinear methods in this literature} \cite{Rootetal2023,Sevgen2024100111,Zilli2024}. {A practical training-time equity intervention should therefore be compatible with these commonly studied and used model classes rather than require a specialized valuation architecture.}

 {Operational adoption by assessment offices nevertheless remains partial and uneven.} \textcite{BidansetRakow2022} identify {communication challenges, limited technical capacity, and budget constraints among perceived barriers to wider adoption}. {Moreover, predictive modeling is only one component of the assessment workflow; the use of ML does not replace sales verification, out-of-time evaluation, ratio studies, model interpretation, or administrative review} \cite{IAAO2020SalesVerification,IAAOAITaskForce2022}. {The} \textcite{IAAOAITaskForce2022} emphasizes explainability, defensibility, administrative feasibility, and integration into existing workflows as central conditions for operational use of advanced ML models. {These considerations favor equity-aware methods that can be incorporated into existing model-training, validation, and review pipelines with minimal disruption.}

\paragraph{Vertical equity and regressivity.}
Equity has long been a central objective in mass appraisal practice and standards. The assessor literature distinguishes \emph{horizontal equity}, which concerns the uniform treatment of similar properties, from \emph{vertical equity}, which concerns systematic differences in assessment outcomes across property-value levels \cite{IAAO2013Ratio,cornia2006property,McMillenSingh2023}. A common form of vertical inequity is \emph{regressivity}: assessment-to-sale ratios tend to decline with property value, so lower-valued properties tend to be relatively overassessed and higher-valued properties tend to be relatively underassessed. {Ratio studies evaluate assessment uniformity using complementary diagnostics. For example, the Coefficient of Dispersion (COD) summarizes dispersion in assessment ratios, whereas the Price-Related Differential (PRD) and Price-Related Bias (PRB) are commonly used to diagnose price-related vertical inequity} \cite{IAAO2013Ratio,IAAO2023VerticalReview}.

% \textcolor{blue}{\st{From the ratio studies measures definitions, one can note that no single statistic fully characterizes vertical equity. Aggregate results of these metrics can conceal bias within geographic or property subgroups, and some regression-based diagnostics can behave poorly, in terms of equity quantification, when estimated values are generated by regression models.}}

{No single diagnostic fully characterizes vertical equity. Methodological studies show that vertical-equity measures can differ materially in their statistical behavior, and that some regression-based diagnostics can be biased when assessed values are themselves generated by regression-based valuation models} \cite{IAAO2023VerticalReview,McMillenSingh2023}. {Accordingly, the assessor literature recommends examining multiple complementary diagnostics rather than relying on a single global measure.}  {Empirical studies further document substantial regressivity and consequential distributional disparities across jurisdictions} \cite{Berry2021,AvenancioLeonHoward2022}. {Recent large-scale evidence also shows that predictive accuracy and assessment equity need not be intrinsically opposed: improvements in the information available to valuation models can, in many settings, improve both simultaneously \cite{SmithEtAl2026}.}

{A smaller literature moves beyond aggregate diagnosis toward model specification, localized validation, and explicit correction intended to improve assessment equity.} Proposed approaches include spatial and geographically weighted methods, stratified procedures, ratio-study-based model validation, and fairness-aware valuation models \cite{Quintos2014,BidansetEtAl2019,BidansetEtAl2025,HermansEtAl2022,HermansEtAl2025,CandoganHanLu2023} \todo{Do we need to sparsify or make more granular comments here? I feel I aggregated all the literature in one chunk}. {Most closely related to our work, \textcite{CandoganHanLu2023} develop a \emph{K-segment} property-valuation model that explicitly targets regressivity and studies the associated accuracy--fairness relationship. Our approach addresses the same assessor-side objective through a different mechanism: rather than introducing a segmented valuation architecture, we incorporate a dependence-based penalty directly into the training objective of existing regression learners.} {The method developed here therefore contributes to this strand by directly targeting a global price-related dependence associated with regressivity during training, while retaining standard ratio studies and subgroup review for out-of-sample evaluation.}

\paragraph{Fair regression and dependence-based constrained learning.}
Although much of the fairness-in-ML literature initially focused on classification, \emph{fair regression} studies how fairness requirements can be incorporated when predictions are continuous, {as in property valuation}. Existing approaches include constrained or regularized empirical risk minimization and post-processing methods. {Different works consider fairness notions such as statistical parity, bounded group loss, mean parity, or distributional parity. Methodologically, these requirements can be imposed through constrained or regularized empirical risk minimization or through post-processing and recalibration procedures} \cite{AgarwalDudikWu2019,FitzsimonsEtAl2019,ChzhenEtAl2020Wasserstein,ChzhenEtAl2020Recalibration,WeiEtAl2023}. {The broader methodological principle is that predictive risk can be optimized subject to, or jointly with, an application-specific fairness criterion. Varying the strength of this requirement produces a family of solutions that makes the relationship between predictive performance and the targeted fairness criterion explicit, without requiring that improvements in one necessarily worsen the other.}

{A particularly relevant strand expresses fairness requirements as restrictions on statistical dependence between predictions and a sensitive variable.} {Prior work has quantified or controlled such dependence using correlation and coefficients of determination \cite{KomiyamaEtAl2018}, kernel-based measures such as the Hilbert--Schmidt Independence Criterion (HSIC) \cite{PerezSuayEtAl2017,LiEtAl2022Fairness}, and maximal-correlation criteria \cite{MaryEtAl2019,lee2022maximal}. Related regularization approaches also control the contribution of sensitive attributes to regression predictions \cite{ScutariPaneroProissl2022}.} {These approaches are particularly relevant when the variable of concern is continuous and can often be incorporated directly into regression objectives. Our application, however, differs from conventional group-fairness settings. Sale price, our observed proxy for market value, is the prediction target rather than a protected attribute, and our objective is not demographic parity. We instead adapt the general dependence-control principle to a different structural relationship: dependence between valuation errors and property value.} {The assessor literature therefore supplies the domain-specific target---price-related regressivity---while the fair-regression literature supplies the general dependence-control mechanism. The proposed method connects the two by translating value-related dependence in valuation errors into a trainable regularization objective.}










\section{Preliminaries: The CCAO Baseline and Empirical Motivation}
\label{sec:setting}

% \textcolor{blue}{\sout{This section fixes the CCAO application setting, prediction notation, and evaluation criteria used throughout the paper. We first describe the residential valuation workflow and baseline model classes, then distinguish predictive performance from assessor-facing ratio-study diagnostics, and finally present the controlled baseline comparison that motivates the proposed correction.}}

{This section fixes the CCAO application setting, prediction notation, and evaluation framework used throughout the paper. We first describe the residential valuation workflow and baseline model classes. We then distinguish predictive performance, assessor-facing measures of assessment level and ratio dispersion, complementary vertical-equity diagnostics, and supplemental checks for nonlinear or localized value-related patterns. Finally, we present the controlled baseline comparison that motivates the proposed correction.}


\subsection{The CCAO Residential Valuation Workflow}
\label{subsec:ccao_workflow}

The CCAO values a large and heterogeneous residential inventory on a triennial reassessment cycle \cite{CCAOAssessmentCalendar2026}. Its residential modeling evolved from township-level linear-regression models, with linear-regression and gradient-boosting alternatives used in 2019--2020, to a county-wide LightGBM model beginning in 2021 \cite{CCAOModelResAVM2026}. LightGBM is an implementation of gradient-boosted decision trees \cite{lightgbm}.

The CCAO's current published residential automated valuation model (AVM) uses verified sales and a broad set of property, spatial, neighborhood, market, and contextual predictors to estimate residential market values \cite{CCAOModelResAVM2026}. The prediction model is embedded in a broader operational workflow comprising data preparation, model training, valuation, evaluation, interpretation, and final review. Candidate valuations are evaluated using both predictive-performance measures and assessor-specific ratio-study statistics across multiple geographic areas and property classes.

% \textcolor{blue}{\sout{For the empirical analysis, we compare a linear-regression benchmark motivated by CCAO's earlier models with standard LightGBM, representing its current model class, while holding fixed the prediction target, underlying feature information, development sample, temporal design, and evaluation procedures.}}

{For the empirical analysis, we compare a linear-regression benchmark motivated by CCAO's earlier models with unpenalized LightGBM, representing its current model class, while holding fixed the prediction target, underlying feature information, development sample, temporal design, and evaluation procedures.}
The comparison is therefore a controlled comparison of model classes rather than a reconstruction of official historical assessments. The exact sample construction, transaction filters, and temporal validation procedures are described in Section~\ref{subsec:ccao_design}.


\subsection{Prediction Setting and Baseline Models}
\label{subsec:baseline_models}

Let $\mathcal{I}$ denote the set of residential sales retained after the CCAO sales-validation and transaction-eligibility procedures described in Section~\ref{subsec:ccao_design}. For each $i\in\mathcal{I}$, we observe
\[
(x_i,P_i,t_i),
\]
where $x_i$ contains property and parcel characteristics; tract-level socioeconomic and demographic variables from the U.S. Census Bureau's American Community Survey (ACS); and geographic, neighborhood, environmental, accessibility, market, and temporal features.

% \textcolor{blue}{\sout{The variable $P_i>0$ is the verified sale price and serves as the transaction-based benchmark for market value, while $t_i$ is the sale date used to construct chronological samples and time-related predictors.}}

{The variable $P_i>0$ is the verified sale price and is treated as the observed transaction-based proxy for the property's market value, while $t_i$ is the sale date used to construct chronological samples and time-related predictors \todo{decide if the $t$ represents any predictor, since they come basically from the same original feature} \cite{CCAOModelResAVM2026,IAAO2026ExposureRatio}. Market value itself is latent: a verified sale therefore provides an observed market-value indicator rather than a directly observed market value.}
The prediction problem is to estimate market value from the information available in $x_i$.

Since residential sale prices are positive and highly right-skewed, the baseline models are trained on the log-price scale. Define
\begin{equation}
\label{eq:prediction_scale}
y_i=\log P_i,
\qquad
\widehat{y}_i=f(x_i),
\qquad
e_i(f)=\widehat{y}_i-y_i,
\end{equation}
where $f\in\mathcal{H}$ is a candidate prediction function. The corresponding prediction on the original price scale is
\begin{equation}
\label{eq:price_prediction}
\widehat{P}_i=\exp\bigl(f(x_i)\bigr).
\end{equation}

% \textcolor{blue}{\sout{The log-scale is the estimation scale, while $\widehat{P}_i$ is the market-value prediction used for price-scale predictive evaluation and ratio-study diagnostics. Because both $\widehat{P}_i$ and $P_i$ are expressed on the market-value scale, define}}

{The log-price scale is the estimation scale, whereas $\widehat{P}_i$ is the directly retransformed valuation used for price-scale predictive evaluation and ratio-study diagnostics. Throughout the analysis, $\widehat{P}_i$ is defined operationally by~\eqref{eq:price_prediction}. Note that the direct exponentiation should not be interpreted as a claim that $\widehat{P}_i$ is a bias-corrected estimate of the expected value of property $i$ given $x_i$. We define the corresponding \emph{valuation-to-sale ratio} as}
\begin{equation}
\label{eq:assessment_ratio_training}
r_i(f)
=
\frac{\widehat{P}_i}{P_i}
=
\exp\bigl(e_i(f)\bigr).
\end{equation}

% \textcolor{blue}{\sout{For brevity and consistency with ratio-study terminology, we refer to $r_i(f)$ as the assessment ratio throughout the paper. Its numerator is the model-estimated market value, not Cook County's statutory assessed value after downstream adjustments.}}

{For brevity and consistency with ratio-study terminology, we subsequently refer to $r_i(f)$ as the \emph{assessment ratio} \todo{re-think about this name. The IAAO uses ratios-to-something, and assessments can be confusing given  the adjustments applied to the prediction, to get the assessment value.}. Its numerator is the model-estimated market value, not Cook County's statutory assessed value after downstream adjustments.}
When a fitted model is fixed, we suppress the dependence on $f$ and write simply $r_i$ and $e_i$.

The log-scale residual and the assessment ratio are therefore linked exactly by
\begin{equation}
\label{eq:log_ratio}
e_i(f)
=
\log \widehat{P}_i-\log P_i
=
\log r_i(f).
\end{equation}
Thus, $e_i(f)>0$ (equivalently, $r_i(f)>1$) corresponds to a prediction above the observed sale price, whereas $e_i(f)<0$ (equivalently, $r_i(f)<1$) corresponds to a prediction below it. Estimation on the log-price scale therefore converts multiplicative price errors into additive residuals, and the training residual is exactly the log assessment ratio.

For the predictive models, we consider two baseline model classes: linear regression and gradient-boosted decision trees. As described above, these choices provide a historically motivated linear benchmark and a representation of CCAO's current model class. Between models, we hold fixed the prediction target, development observations, underlying feature information, temporal design, and evaluation protocol; the models differ in functional form and model-specific estimation procedure.

For a training sample of size $n$, both baselines are fitted to reduce mean squared log-price error, i.e.,
\begin{equation}
\label{eq:baseline_learning}
L(f)
=
\frac{1}{n}\sum_{i=1}^n e_i(f)^2,
\end{equation}
while differing in functional form, estimation procedure, and model-specific complexity control. This loss is the predictive component of the training objective modified in Section~\ref{sec:method}.

For linear regression,
\begin{equation}
\label{eq:linear_model}
f_{\mathrm{lin}}(x)
=
\beta_0+x^\top\beta.
\end{equation}
The hypothesis class is therefore restricted to functions that are linear in the supplied predictors. In the ordinary linear-regression baseline, no additional regressivity term is included in the training objective. This specification provides the controlled linear benchmark motivated by CCAO's earlier linear-regression models. Its linear structure imposes a comparatively rigid relationship between the predictors and log sale price.

% \textcolor{blue}{%
  % \sout{LightGBM, the modeling foundation of CCAO's current published residential AVM}%%
  % \sout{, constructs an additive ensemble of regression trees.}%
% }

{On the other hand, LightGBM, the modeling foundation of CCAO's current published residential AVM \todo{Confirm if I defined AVM or if I should use another term instead}, constructs an additive ensemble of regression trees \cite{CCAOModelResAVM2026,lightgbm}. A gradient-boosted tree predictor can be written as}
\begin{equation}
\label{eq:boosted_tree_model}
f_{\mathrm{gbm}}^{(M)}(x)
=
f^{(0)}(x)
+
\sum_{m=1}^{M}\nu b_m(x),
\end{equation}
where $f^{(0)}$ is an initial prediction, $b_m$ is the regression tree added at boosting iteration $m$, $M$ is the number of boosting iterations, and $\nu>0$ is the learning rate. Unlike the linear model, the tree ensemble can represent nonlinearities, threshold effects, and interactions without requiring them to be specified in advance.

LightGBM fits the ensemble in~\eqref{eq:boosted_tree_model} sequentially. The gradient and Hessian calculations required for the proposed custom objectives are introduced with the penalties in Section~\ref{sec:method}.

The two baselines therefore solve the same prediction problem under a common empirical design but use substantially different model classes. Both baseline objectives minimize squared log-price error and contain no term that directly controls systematic value-related patterns in residuals or assessment ratios. We next distinguish predictive performance from the assessor-facing criteria used to evaluate assessment level, ratio dispersion, and regressivity.


\subsection{Predictive and Assessor-Facing Evaluation Criteria}
\label{subsec:metrics}

% \textcolor{blue}{\sout{The two scales play distinct roles in evaluation. Predictive performance is assessed both on the log-price scale used for estimation and on the original price scale on which market values are reported. The assessor-facing ratio-study diagnostics are computed on the original price scale.}}

{The two scales play distinct roles in evaluation. Predictive performance is assessed both on the log-price scale used for estimation and on the original price scale on which valuations are reported. Assessor-facing ratio-study diagnostics are computed on the original price scale. The resulting measures answer different questions and are therefore interpreted jointly rather than collapsed into a single performance criterion.}

\paragraph{Predictive performance.}
% \textcolor{blue}{\sout{For an evaluation sample $S\subseteq\mathcal{I}$, let $n_S=|S|$, $\bar P_S=n_S^{-1}\sum_{i\in S}P_i$, and $\bar y_S=n_S^{-1}\sum_{i\in S}y_i$. We use three complementary headline measures: price-scale mean absolute error (MAE), price-scale $R^2$, and log-scale root mean squared error (RMSE).}}
{For an evaluation sample $S\subseteq\mathcal{I}$, let $n_S=|S|$ and $\bar P_S=\frac{1}{n_S}\sum_{i\in S}P_i$. We use four complementary headline accuracy-related measures: price-scale mean absolute error (MAE), price-scale mean absolute percentage error (MAPE), price-scale $R^2$, and log-scale root mean squared error (RMSE):}
\begin{align}
\label{eq:predictive_metrics}
\operatorname{MAE}_{P}(S)
&=
\frac{1}{n_S}\sum_{i\in S}\left|\widehat P_i-P_i\right|,
\\
{\operatorname{MAPE}_{P}(S)}
&{=100\times
\frac{1}{n_S}
\sum_{i\in S}
\left|
\frac{\widehat P_i-P_i}{P_i}
\right|
=100\times 
\frac{1}{n_S}\sum_{i\in S}|r_i-1|},
\\
R^2_{P}(S)
&=
1-
\frac{\sum_{i\in S}(\widehat P_i-P_i)^2}
{\sum_{i\in S}(P_i-\bar P_S)^2},
\\
\operatorname{RMSE}_{\log P}(S)
&=
\left(
\frac{1}{n_S}\sum_{i\in S}e_i(f)^2
\right)^{1/2}.
\end{align}
% \textcolor{blue}{\sout{The price-scale MAE gives the average absolute prediction error in dollars, while $R^2_P$ summarizes predictive fit on the market-value scale. The log-scale RMSE directly evaluates the loss scale used for estimation and, by~\eqref{eq:log_ratio}, can equivalently be viewed as the root mean squared log-ratio error. When useful, we also report the supplementary log-scale coefficient of determination $R^2_{\log P}$.}}

{The price-scale MAE gives the average absolute prediction error in dollars, while MAPE gives the average absolute error relative to the observed sale price. In particular, MAPE measures deviation from the economically neutral ratio $r_i=1$ and is therefore distinct from COD below, which measures dispersion around the sample median ratio. The statistic $R^2_P$ summarizes predictive fit on the original price scale. Finally, $\operatorname{RMSE}_{\log P}$ directly evaluates the scale and squared-error loss used for estimation and, by~\eqref{eq:log_ratio}, is equivalently the root mean squared log-ratio error.}
% \textcolor{blue}{\sout{Because $R^2_{\log P}$ and $\operatorname{RMSE}_{\log P}$ are driven by the same log-scale sum of squared errors on a fixed evaluation sample, $R^2_{\log P}$ is treated as supplementary rather than as a separate headline criterion. Percentage errors are also not used as headline measures: for example, on the price scale, the mean absolute percentage error (MAPE) equals $\frac{100}{n_S}\sum_{i\in S}|r_i-1|$ and therefore overlaps directly with the ratio-based evaluation below, while percentage errors of log-prices do not have a stable scale interpretation.}}
% \textcolor{impGreen}{We retain $R^2_{\log P}$ only as a supplementary implementation diagnostic. On a fixed evaluation sample it is driven by the same log-scale sum of squared errors as $\operatorname{RMSE}_{\log P}$ and is therefore not treated as an additional headline criterion.}
\textcolor{impGreen}{\todo{Reconcile the code/output ``OOS R2'' with the definitions in this subsection. Do we use train mean as baseline or the test one, as usually done?}}

\paragraph{Assessment level.}
Recall from~\eqref{eq:assessment_ratio_training} that $r_i=\widehat P_i/P_i$. A ratio above one indicates that the predicted value exceeds the observed sale price, whereas a ratio below one indicates that it falls below the sale price. Define the median, arithmetic mean, and price-weighted mean ratios as
\begin{equation}
\label{eq:level_ratios}
m_S
=
\operatorname{median}_{i\in S}\{r_i\},
\qquad
\bar r_S
=
\frac{1}{n_S}\sum_{i\in S}r_i,
\qquad
\bar r_{W,S}
=
\frac{\sum_{i\in S}\widehat P_i}{\sum_{i\in S}P_i}
=
\frac{\sum_{i\in S}r_iP_i}{\sum_{i\in S}P_i}.
\end{equation}

These quantities summarize aggregate assessment level. Because both numerator and denominator of $r_i$ are expressed on the market-value scale in this analysis, the neutral reference value for the three statistics is a value of one. The median is relatively insensitive to extreme ratios, the arithmetic mean assigns equal weight to each sale, and the weighted mean places greater weight on higher-value properties. Agreement among these measures indicates broad aggregate calibration but does not establish either horizontal or vertical equity.

% \textcolor{impGreen}{Any numerical assessment-level bands used as model-selection guardrails are specified separately in Section~\ref{sec:model_selection} \todo{we erased the model selection section. Update this reference}; they are not treated here as universal substitutes for the neutral target of one.}

\paragraph{Ratio dispersion (Horizontal uniformity).}
Ratio dispersion is an important assessor-facing dimension of horizontal uniformity. We summarize it using the coefficient of dispersion (COD) and coefficient of variation (COV). The COD is
\begin{equation}
\label{eq:cod}
\operatorname{COD}(S)
=
100\times
\frac{1}{n_S}
\sum_{i\in S}
\frac{|r_i-m_S|}{m_S},
\end{equation}
and the COV is
\begin{equation}
\label{eq:cov}
\operatorname{COV}(S)
=
100\times
\frac{s_{r,S}}{\bar r_S},
\qquad
s_{r,S}
=
\left(
\frac{1}{n_S-1}
\sum_{i\in S}(r_i-\bar r_S)^2
\right)^{1/2}.
\end{equation}

% \textcolor{blue}{\sout{The COD measures the average absolute deviation from the median ratio, whereas the COV measures standard deviation relative to the mean. Lower values indicate less aggregate ratio dispersion. Applicable reference ranges depend on property class and prevailing market conditions.}}

{The COD measures the average absolute deviation from the median ratio, whereas the COV measures standard deviation relative to the mean. Lower values indicate less aggregate ratio dispersion, but the reference ranges depend on property class and prevailing market conditions, and unusually low dispersion can itself warrant diagnostic review \cite{IAAO2013Ratio,IAAO2026ExposureRatio}.}
To avoid ambiguity with the covariance penalty introduced later, $\operatorname{COV}$ always denotes the coefficient of variation, whereas statistical covariance is written as $\operatorname{Cov}(\cdot,\cdot)$.

\paragraph{Regressivity (Vertical equity).}
Vertical equity concerns whether assessment ratios vary systematically with property value. A \emph{regressive} pattern arises when lower-value properties tend to have higher assessment ratios than higher-value properties. A \emph{progressive} pattern arises in the opposite direction.

{No single statistic fully characterizes vertical equity: different diagnostics summarize different aspects of the relationship between valuations, ratios, and property value and can therefore disagree \cite{IAAO2023VerticalReview,IAAO2026ExposureRatio}. We consequently use a complementary suite of ratio-based, regression-based, distributional, and grouped diagnostics rather than treating any single statistic as definitive. The May 2026 IAAO exposure draft places primary emphasis on the grouped VEI and identifies PRB, MKI, and PRD as complementary measures. Given that the document is still in draft form, we use its proposed procedures and ranges as descriptive guidance rather than as final adopted compliance standards.}

The \emph{price-related differential} (PRD) is
\begin{equation}
\label{eq:prd}
\operatorname{PRD}(S)
=
\frac{\bar r_S}{\bar r_{W,S}}.
\end{equation}
Note that, since the weighted mean ratio assigns greater weight to higher-value properties, a PRD above one tends to arise when higher-value properties have lower assessment ratios than lower-value properties. The 2013 IAAO standard provides a reference range of $[0.98,1.03]$ for the PRD: values above this range tend to indicate regressivity, whereas values below it tend to indicate progressivity \cite{IAAO2013Ratio}.

{PRD is a simple and widely used ratio-based diagnostic, but it can be particularly sensitive to influential high-value observations and extreme ratios. Therefore, we interpret it jointly with the other vertical-equity measures rather than as a standalone test \cite{IAAO2026ExposureRatio}.}

% \textcolor{blue}{\sout{We also report the price-related bias (PRB). Following the IAAO PRB methodology, define the equal-weight market-value proxy}}

{The following regression- and group-based vertical-equity diagnostics require a proxy for the property's value level. Ranking properties only by sale price can mechanically bias ratio-based vertical-equity tests toward finding regressivity because sale price also appears in the denominator of $r_i$, while using valuation alone can induce bias in the opposite direction. Following IAAO guidance, we therefore use the equal-weight market-value proxy}
\begin{equation}
\label{eq:market_value_proxy}
V_i^{\mathrm{proxy}}
=
\frac{1}{2}
\left(
P_i+\frac{\widehat P_i}{m_S}
\right),
\end{equation}
{which combines the verified sale price with the median-normalized predicted value to reduce this source of proxy-induced bias \cite{IAAO2013Ratio,IAAO2026ExposureRatio}.}

{We first report the \textit{price-related bias} (PRB).}
Let $\widetilde V_S^{\mathrm{proxy}}$ denote the median of $V_i^{\mathrm{proxy}}$ over $S$. The PRB is the slope in the simple linear regression
\begin{equation}
\label{eq:prb}
\frac{r_i-m_S}{m_S}
=
a_S
+
\operatorname{PRB}(S)\cdot
\log_2\left(
\frac{V_i^{\mathrm{proxy}}}
{\widetilde V_S^{\mathrm{proxy}}}
\right)
+
\varepsilon_i.
\end{equation}
\todo{I think we can improve the notation here for the slope PRB}
Negative PRB values indicate that assessment ratios tend to decline as property value increases and are therefore associated with regressivity. Values between $-0.05$ and $0.05$ are conventionally preferred \cite{IAAO2013Ratio}.

{As a complementary distributional diagnostic, we report the \emph{Modified Kakwani Index} (MKI). Let $\pi(1),\allowbreak\ldots,\pi(n_S)$ index the observations in $S$ in nondecreasing order of sale price, and let
\[
q_j=\frac{j-\tfrac12}{n_S},
\qquad
\overline{\widehat P}_S
=
\frac{1}{n_S}\sum_{i\in S}\widehat P_i.
\]
Define the \emph{concentration index} of predicted values when observations are ordered by sale price and the corresponding \emph{Gini coefficient} of sale prices as
\begin{equation}
\label{eq:mki_components}
C_{\widehat P\mid P}(S)
=
\frac{2}{n_S\overline{\widehat P}_S}
\sum_{j=1}^{n_S}
\widehat P_{\pi(j)}q_j
-1,
\qquad 
G_P(S)
=
\frac{2}{n_S\bar P_S}
\sum_{j=1}^{n_S}
P_{\pi(j)}q_j
-1,
\end{equation}
respectively. Then, the MKI is
\begin{equation}
\label{eq:mki}
\operatorname{MKI}(S)
=
\frac{C_{\widehat P\mid P}(S)}{G_P(S)}.
\end{equation}
When there is sufficient price variation, an MKI value of one indicates vertical neutrality, whereas values below one indicate a regressive direction, and values above one indicate a progressive direction \cite{IAAO2026ExposureRatio}. Because MKI orders observations using sale price, it can itself be affected by errors in the market-value proxy, which is another reason to interpret it as one component of the broader diagnostic suite rather than as a definitive test.}
\todo{Verify that the implementation of MKI uses the price-ordered concentration index in~\eqref{eq:mki_components}, and not the simplified ratio $\operatorname{Gini}(\widehat P)/\operatorname{Gini}(P)$, which is not the MKI definition used here.}

% \textcolor{impGreen}{}

% \textcolor{blue}{\sout{Finally, we report the Vertical Equity Indicator (VEI), a high-versus-low value-group comparison proposed in the May 2026 IAAO exposure draft.}}

Finally, we report the \emph{Vertical Equity Indicator} (VEI), the quantile-grouped vertical-equity measure emphasized in the May 2026 IAAO exposure draft \cite{IAAO2026ExposureRatio}.
The draft varies the number of percentile groups with sample size: for the large evaluation samples used in the main countywide analysis ($n_S\geq 501$), observations are divided into deciles according to $V_i^{\mathrm{proxy}}$. Let $S_{q_1}$ and $S_{q_{10}}$ denote the lowest and highest deciles, and let $m_{S_{q_1}}$ and $m_{S_{q_{10}}}$ denote their median assessment ratios. The VEI is
\begin{equation}
\label{eq:vei}
\operatorname{VEI}(S)
=
100\times
\frac{m_{S_{q_{10}}}-m_{S_{q_1}}}{m_S}.
\end{equation}

A negative VEI indicates that higher-value properties have lower median assessment ratios than lower-value properties. The exposure draft proposes a reference band of $[-10\%,10\%]$, together with confidence-interval-based testing for point estimates outside that band.

% \textcolor{blue}{\sout{Because the May 2026 document remains an exposure draft, we use VEI and the proposed band as supplementary descriptive guidance rather than as a final adopted compliance standard.}}
{Again, since the May 2026 document remains an exposure draft, we use VEI and the proposed band as descriptive guidance rather than as a final adopted compliance standard. Moreover, note that the first-versus-last quantile comparison in~\eqref{eq:vei} does not necessarily reveal non-monotone patterns within the median ratios distribution. Hence, we also report the median assessment ratio and its 90\% confidence interval for every VEI percentile group, together with the overall median ratio.} 

% For smaller stratified samples, the number of groups follows the corresponding sample-size guidance in the exposure draft rather than imposing deciles mechanically \cite{IAAO2026ExposureRatio}.}

% \textcolor{blue}{\sout{The previous summary table, which omitted MAPE and MKI and summarized vertical equity only through PRD, PRB, and VEI, is replaced by the revised table below.}}


\begin{table}[!ht]
\centering
\footnotesize
\setlength{\tabcolsep}{3.2pt}
\renewcommand{\arraystretch}{1.10}

% Indentation for individual measures within each family
\newcommand{\metriclabel}[1]{\hspace*{0.9em}#1}

\caption{Predictive and assessor-facing evaluation measures.}
\label{tab:assessment_metrics_summary}

\begin{tabularx}{\textwidth}{
@{}
>{\raggedright\arraybackslash}p{3.05cm}
>{\centering\arraybackslash}p{0.80cm}
>{\centering\arraybackslash}p{1.70cm}
>{\centering\arraybackslash}p{1.40cm}
>{\raggedright\arraybackslash}X
@{}
}
\toprule
\textbf{Measure}
& \textbf{\shortstack{Ideal\\value}}
& \textbf{\shortstack{Reference\\range}}
& \textbf{\shortstack{Ideal\\direction}}
& \textbf{Interpretation} \\
\midrule

% \multicolumn{5}{l}{\textit{Predictive Performance}}\\
\textit{Predictive performance}
& & & & \\
\cmidrule(r){1-1}

\metriclabel{$R^2_P$}
& $1$
& --
& Higer $\uparrow$
& Original price-scale goodness of fit. \\

\metriclabel{$\operatorname{MAE}_P$}
& $0$
& --
& Lower $\downarrow$
& Mean absolute prediction error in dollars. \\

\metriclabel{$\operatorname{MAPE}_P$}
& $0$
& --
& Lower $\downarrow$
& Mean absolute prediction error relative to sale price. \\

\metriclabel{$\operatorname{RMSE}_{\log P}$}
& $0$
& --
& Lower $\downarrow$
& RMSE on the log-price estimation scale. \\


\midrule
\addlinespace[5pt]
\textit{Assessment level}
& & & & \\
\cmidrule(r){1-1}

\metriclabel{Median ratio $m_S$}
& $1$
& --
& $\rightarrow 1$
& Robust aggregate assessment level. \\

\metriclabel{Mean ratio $\bar r_S$}
& $1$
& --
& $\rightarrow 1$
& Unweighted aggregate assessment level. \\

\metriclabel{Weighted mean $\bar r_{W,S}$}
& $1$
& --
& $\rightarrow 1$
& Value-weighted assessment level. \\

\midrule
\addlinespace[6pt]
% \shortstack[l]{
\textit{Horizontal equity}%\\[-1pt]
% \textit{\scriptsize Horizontal uniformity}
% }
& & & & \\
\cmidrule(r){1-1}

\metriclabel{COD}
& --
& $[5.0,15.0]$\textsuperscript{$\dagger$}
& $\rightarrow$ range
& Absolute dispersion relative to the median ratio. \\

\metriclabel{COV}
& --
& --
& Lower $\downarrow$
& Relative dispersion relative to the mean ratio. \\

\midrule
\addlinespace[5pt]
% \shortstack[l]{
\textit{Vertical equity}%\\[-1pt]
% \textit{\scriptsize Regressivity}
% }
& & & & \\
\cmidrule(r){1-1}

\metriclabel{PRD}
& $1$
& $[0.98,1.03]$
& $\rightarrow 1$
& Mean-to-weighted-mean ratio comparison. \\

\metriclabel{PRB}
& $0$
& $[-0.05,0.05]$
& $\rightarrow 0$
& Regression-based value-related ratio trend. \\

\metriclabel{MKI}
& $1$
& --
& $\rightarrow 1$
& Distributional concentration relative to sale prices. \\

\metriclabel{VEI}
& $0$
& $[-10,10]$\textsuperscript{$\dagger$}
& $\rightarrow 0$
& High-versus-low value-group median-ratio difference. \\

\bottomrule
\end{tabularx}

\vspace{1mm}
\begin{minipage}{\textwidth}
\scriptsize
\emph{Notes.}
``--'' denotes that no universal ideal or reference range is imposed.
Arrows indicate the preferred direction of improvement.
PRD and PRB reference ranges follow IAAO guidance \cite{IAAO2013Ratio}.
\textsuperscript{$\dagger$}Reference ranges are from the May 2026
IAAO exposure draft: applicable COD ranges depend on type of property and market
profile \cite{IAAO2026ExposureRatio}.
\end{minipage}

\end{table}



% \begin{table}[!ht]
% % \color{impGreen}
% \centering
% \caption{Summary of predictive and assessor-facing evaluation measures.}
% \label{tab:assessment_metrics_summary}
% \begin{tabular}{lcp{8.0cm}}
% \toprule
% \textbf{Measure} & \textbf{Reference/Direction} & \textbf{Interpretation} \\
% \midrule
% \multicolumn{3}{l}{\textit{Predictive Performance}} \\
% $R^2_P$ & Higher is better & Predictive fit on the original price scale. \\
% $\operatorname{MAE}_P$ & Lower is better & Average absolute prediction error in dollars. \\
% $\operatorname{MAPE}_P$ & Lower is better & Average absolute prediction error relative to sale price. \\
% $\operatorname{RMSE}_{\log P}$ & Lower is better & Prediction error on the log-price scale used for estimation. \\
% \midrule
% \multicolumn{3}{l}{\textit{Assessment Level}} \\
% Median ($m_S$) & $1$ & Outlier-robust measure of aggregate assessment level. \\
% Mean ($\bar r_S$) & $1$ & Arithmetic mean assessment ratio. \\
% Weighted mean ($\bar r_{W,S}$) & $1$ & Value-weighted mean ratio; denominator in PRD. \\
% \midrule
% \multicolumn{3}{l}{\textit{Horizontal Uniformity (Ratio Dispersion)}} \\
% COD & $[5.0,15.0]$\textsuperscript{*}
% & Average percentage deviation from the median ratio; applicable ranges depend on property and market profile. \\
% COV & Lower is better
% & Standard deviation relative to the mean ratio, reported in percentage points. \\
% \midrule
% \multicolumn{3}{l}{\textit{Vertical Equity (Regressivity)}} \\
% VEI & $[-10\%,10\%]$\textsuperscript{$\dagger$}
% & Grouped high-versus-low value comparison; negative values indicate a regressive direction. \\
% PRB & $[-0.05,0.05]$
% & Regression-based diagnostic; negative values outside the range indicate a regressive direction. \\
% MKI & $1$
% & Distributional concentration diagnostic; values below one indicate a regressive direction. \\
% PRD & $[0.98,1.03]$
% & Ratio-based diagnostic; values above the range tend to indicate a regressive direction. \\
% \bottomrule
% \end{tabular}

% \vspace{1mm}
% \begin{minipage}{0.96\textwidth}
% \footnotesize
% \textsuperscript{*}The May 2026 IAAO exposure draft gives $5.0$--$15.0$ as a general range for improved residential property; applicable ranges depend on market and property profile, and unusually low dispersion can itself warrant review \cite{IAAO2026ExposureRatio}.\\
% \textsuperscript{$\dagger$}VEI, its proposed reference band, and the corresponding percentile-group procedures are from the nonfinal May 2026 IAAO exposure draft \cite{IAAO2026ExposureRatio}. At the time of writing, the exposure draft is not a final adopted standard; the 2013 Standard on Ratio Studies remains the operative adopted reference where applicable \cite{IAAO2013Ratio}.
% \end{minipage}
% \end{table}

% \textcolor{blue}{\sout{Table~\ref{tab:assessment_metrics_summary} summarizes the principal evaluation criteria used throughout the paper. Predictive metrics quantify valuation error, assessment-level statistics summarize aggregate calibration, COD and COV summarize ratio dispersion, and PRD, PRB, and VEI diagnose price-related patterns. These dimensions should be interpreted jointly: strong predictive performance or low aggregate dispersion does not by itself imply neutral vertical-equity behavior.}}

{Table~\ref{tab:assessment_metrics_summary} summarizes the principal predictive and assessor-facing criteria used throughout the paper. Predictive metrics quantify different aspects of valuation error, assessment-level statistics summarize aggregate calibration, COD and COV summarize ratio dispersion, and VEI, PRB, MKI, and PRD provide complementary grouped, regression-based, distributional, and ratio-based views of vertical equity. Strong predictive performance or low aggregate dispersion does not by itself imply neutral vertical-equity behavior.}

{
\color{impGreen}
\paragraph{Supplemental diagnostics.}
The assessor-facing vertical-equity measures are complemented by diagnostics
that examine the relationship targeted by the proposed correction more
directly. First, we report the slope from regressing log ratios on log sale
prices:
\begin{equation}
\label{eq:log_ratio_slope}
\beta_{\log}(S)
=
\frac{
\widehat{\operatorname{Cov}}_S(e,y)
}{
\widehat{\operatorname{Var}}_S(y)
}
=
\frac{
\widehat{\operatorname{Cov}}_S(\log r,\log P)
}{
\widehat{\operatorname{Var}}_S(\log P)
}.
\end{equation}
A value of zero indicates no first-order log-ratio--log-sale-price trend.
Negative values indicate that ratios tend to decline with sale price, the
regressive direction, whereas positive values indicate the progressive
direction. 
% Because $\widehat{\operatorname{Var}}_S(y)$ is fixed for a given
% evaluation sample, $\beta_{\log}(S)$ is the slope corresponding directly to
% the covariance targeted during training and contains the same signed
% first-order information when models are compared on the same sample.
% Because sale price is an imperfect proxy for latent market value and also
% appears in the denominator of the ratio, we treat this slope as a
% mechanism-oriented supplemental diagnostic rather than as a standalone
% definition of vertical equity.
Second, we report the \emph{distance correlation}
\begin{equation}
\label{eq:dcor_diagnostic}
\operatorname{dCor}(e,y)
=
\operatorname{dCor}(\log r,\log P),
\end{equation}
as an omnibus diagnostic of dependence between log residuals and log sale
prices \cite{SzekelyRizzoBakirov2007}. Unlike the signed first-order slope,
distance correlation can detect nonlinear and other non-first-order
dependence but does not identify a regressive or progressive direction.
Moreover, distributional independence is stronger than vertical neutrality
as conventionally evaluated in ratio studies. We therefore use
$\operatorname{dCor}(e,y)$ only as a supplemental diagnostic of remaining
value-related dependence, not as an assessor-facing measure of vertical
equity.
}

% \paragraph{\textcolor{impGreen}{Supplemental diagnostics.}} \todo{Note: I think we need this or an equivalent metric to also account for the non-linear distributional independence. We already talked about this in the past, but I am making it explicit here. IAAO metrics are not enough IMO. GOing more granular to locaton or other sensitive attributes buys a lot of credibility too.
% We can sill move it to Appendix if needed...}
% \textcolor{impGreen}{The assessor-facing statistics are complemented by three checks designed around failure modes of the proposed first-order correction. First, we examine the full VEI percentile-group profile described above rather than only its first-versus-last summary. Second, we report the \emph{distance correlation}.}
% \begin{equation}
% \label{eq:dcor_diagnostic}
% \textcolor{impGreen}{
% \operatorname{dCor}\bigl(e,y\bigr)
% =
% \operatorname{dCor}\bigl(\log r,\log P\bigr),
% }
% \end{equation}
% \textcolor{impGreen}{as an omnibus diagnostic of residual value dependence \cite{SzekelyRizzoBakirov2007}. Under standard moment conditions, population distance correlation equals zero if and only if the variables are distributionally independent ($r \perp P$). Note that independence is stronger than vertical neutrality as conventionally evaluated in ratio studies. For example, distance correlation can respond to price-related changes in residual dispersion even when group medians are flat. Hence, we use $\operatorname{dCor}$ only to test whether a (linear) covariance-based correction induces nonlinear dependencies, not as an assessor compliance metric. Third, where sample size and value variation are sufficient, we recompute the principal vertical-equity diagnostics previously defined within prespecified geographic and property strata rather than relying exclusively on countywide averages \cite{IAAO2026ExposureRatio}.}

% \textcolor{blue}{\sout{Two qualifications are important for interpretation. First, throughout the empirical model comparisons we use these point estimates as descriptive diagnostics rather than as formal determinations of compliance with IAAO standards; formal compliance testing can require additional uncertainty analysis. Second, countywide statistics can conceal geographic, property-class, temporal, or nonlinear heterogeneity. We therefore evaluate models using several complementary statistics and ratio plots rather than drawing conclusions from any single measure.}}

{Three benchmark principles govern the empirical use of these measures. First, point estimates are treated as descriptive model diagnostics rather than as formal determinations of compliance with IAAO standards: uncertainty is reported where it is material to substantive comparisons, including the confidence intervals accompanying the VEI percentile-group profile \cite{IAAO2013Ratio,IAAO2026ExposureRatio}. Second, countywide statistics are complemented by prespecified geographic and property-stratum analyses when the available sample and value variation support meaningful comparison. Third, during model development the evaluation statistics are computed on chronological validation folds before the aggregation used for model selection. The full evaluation suite is not itself imposed as a collection of constraints over the model, but as a general diagnostic of its estimates.

% : Section~\ref{sec:model_selection} specifies the subset used for model selection, while the remaining diagnostics provide independent validation against improving only the criteria directly used for selection.}




\subsection{Baseline Comparison and Empirical Motivation}
\label{subsec:baseline_tension}

% -------------------------------------------------------------------------
% Comparison setup
% -------------------------------------------------------------------------

\textcolor{blue}{\sout{In this subsection we compare the historically motivated linear regression benchmark with the current model class used by the CCAO: unpenalized LightGBM.}}

\textcolor{impGreen}{
We compare the historically motivated linear-regression benchmark with
unpenalized LightGBM, representing CCAO's current model class.
}

\textcolor{blue}{\sout{To facilitate a controlled comparison of model classes, we reestimate both models using the same outcome definition, development sample, underlying feature information, and surrounding data-processing and evaluation procedures. The comparison is therefore not a direct historical comparison of official assessments produced in different years. Instead, it evaluates the two model classes under a common empirical design.}}

\textcolor{impGreen}{
Both models are estimated and evaluated under the same prediction target,
underlying feature information, sample-construction rules, temporal-validation
framework, and evaluation procedure. The comparison is therefore a controlled
comparison of model classes rather than a reconstruction of official
assessments produced under different historical workflows.
Section~\ref{subsec:ccao_design} provides the complete empirical design.
}

\textcolor{blue}{\sout{The final empirical design uses sales from 2016--2023 for model development, including all model and hyperparameter selection through chronological validation. Candidate configurations are evaluated on chronological development folds before any across-fold aggregation. The held-out 2024 sample provides the primary out-of-sample baseline comparison and is not used for model or hyperparameter selection. The 2025 assessment-year sample is reserved for the separate forward evaluation reported later and is likewise excluded from selection. Section~\ref{subsec:ccao_design} gives the complete temporal design and selection protocol. The unpenalized LightGBM benchmark is tuned entirely within this development framework before the regressivity penalties are evaluated.}}

\textcolor{impGreen}{
\todo{Once the final temporal design is locked, add at most one sentence here
identifying the held-out sample used for the baseline comparison. Keep the
complete development, validation, and forward-evaluation chronology in
Section~\ref{subsec:ccao_design}.}
}


% -------------------------------------------------------------------------
% Baseline table
% -------------------------------------------------------------------------

% \textcolor{blue}{\sout{Insert the final 2024 baseline comparison. The headline table should report $R^2_P$, $\operatorname{MAE}_P$, $\operatorname{MAPE}_P$, $\operatorname{RMSE}_{\log P}$, the median assessment ratio, COD, VEI, PRB, MKI, and PRD. Keep mean ratio, weighted mean ratio, COV, $R^2_{\log P}$, distance correlation, full VEI profiles, and stratified results in the corresponding supplementary/robustness panels unless they are necessary to explain a headline result. Report the 2025 forward evaluation separately in the empirical-results section.}}

\textcolor{blue}{\sout{\textit{Old placeholder table removed: a horizontally resized two-row table reported an obsolete test-set snapshot and omitted several principal measures introduced in Section~\ref{subsec:metrics}.}}}

{\color{impGreen}

\todo{Populate Table~\ref{tab:ccao_baseline_results} after the final empirical
design is locked. All entries must be computed from the same held-out
predictions and the same final metric implementation.}

\begin{table}[!ht]
\centering
\scriptsize
\setlength{\tabcolsep}{2.2pt}
\renewcommand{\arraystretch}{1.06}

\newcommand{\baselinemetric}[1]{\hspace*{0.6em}#1}

\caption{Baseline comparison on the held-out test sample and the 2025
assessment-year sample.}
\label{tab:ccao_baseline_results}

\begin{tabularx}{\textwidth}{
@{}
>{\raggedright\arraybackslash}p{2.80cm}
>{\centering\arraybackslash}X
>{\centering\arraybackslash}X
>{\centering\arraybackslash}X
>{\centering\arraybackslash}X
@{}
}
\toprule
\textbf{Measure}
& \multicolumn{2}{c}{\textbf{Test}}
& \multicolumn{2}{c}{\textbf{2025 assessment}} \\
\cmidrule(lr){2-3}\cmidrule(l){4-5}
&
\textbf{Linear}
& \textbf{\shortstack{Light\\GBM}}
& \textbf{Linear}
& \textbf{\shortstack{Light\\GBM}} \\
\midrule

\multicolumn{5}{@{}l}{\textbf{Prediction}} \\

\baselinemetric{$R^2_P$}
& 0.799 & \textbf{0.882} & 0.799 & \textbf{0.891} \\

\baselinemetric{$\operatorname{MAE}_P$}
& \$90,101 & \textbf{\$79,118} & \$99,368 & \textbf{\$83,622} \\

\baselinemetric{$\operatorname{MAPE}_P$}
& 24.1\% & \textbf{21.8\%} & 24.9\% & \textbf{21.3\%} \\

\baselinemetric{$\operatorname{RMSE}_{\log P}$}
& 0.322 & \textbf{0.294} & 0.313 & \textbf{0.285} \\


\addlinespace[5pt]
\multicolumn{5}{@{}l}{\textbf{Assessment level}} \\

\baselinemetric{Median ratio $m_S$}
& \textbf{0.969} & 0.918 & \textbf{1.020} & 0.931 \\

\baselinemetric{Mean ratio $\bar r_S$}
& 1.020 & \textbf{0.983} & 1.081 & \textbf{0.996} \\

\baselinemetric{Weighted mean $\bar r_{W,S}$}
& \textbf{0.979} & 0.913 & \textbf{1.027} & 0.918 \\


\addlinespace[5pt]
\multicolumn{5}{@{}l}{\textbf{Horizontal equity}} \\

\baselinemetric{COD}
& 24.7\% & \textbf{22.2\%} & 24.3\% & \textbf{21.7\%} \\

\baselinemetric{COV}
& 45.2\% & \textbf{40.4\%} & 42.1\% & \textbf{37.7\%} \\


\addlinespace[5pt]
\multicolumn{5}{@{}l}{\textbf{Vertical equity}} \\

\baselinemetric{PRD}
& \textbf{1.042} & 1.077 & \textbf{1.052} & 1.085 \\

\baselinemetric{PRB}
& \textbf{-0.016} & -0.103 & \textbf{-0.029} & -0.114 \\

\baselinemetric{MKI}
& \textbf{0.975} & 0.910 & \textbf{0.954} & 0.897 \\

\baselinemetric{VEI}
& \textbf{-11.6\%} & -29.1\% & \textbf{-17.4\%} & -33.1\% \\

\addlinespace[5pt]
\multicolumn{5}{@{}l}{\textbf{Supplemental diagnostic}} \\

\baselinemetric{$\beta_{\log}$}
& \textbf{-0.092} & -0.163 & \textbf{-0.109} & -0.174 \\

\baselinemetric{$\operatorname{dCor}(e,y)$}
& \textbf{0.250} & 0.408 & \textbf{0.269} & 0.435 \\

\bottomrule
\end{tabularx}

\vspace{1mm}
\begin{minipage}{\textwidth}
\scriptsize
\emph{Notes.}
Preferred values, directions, and applicable reference ranges are defined in
Table~\ref{tab:assessment_metrics_summary}. MAPE, COD, COV, and VEI are reported
in percentage points and MAE in dollars. Under the current protocol, test metrics use models fit on the
344,610-sale development pool, while 2025 metrics use production models refit on all
382,900 eligible 2016--2024 sales. $\beta_{\log}$ is defined in Eq.~\eqref{eq:log_ratio_slope}. Distance correlation is computed between the log
residual $e$ and log sale price $y$.
\end{minipage}
\end{table}

}


% -------------------------------------------------------------------------
% Baseline visualization
% -------------------------------------------------------------------------

% \textcolor{blue}{\sout{\textit{Old placeholder figure removed: a single-model scatterplot of assessment ratios against log sale prices did not directly visualize the linear-regression versus LightGBM comparison or the value-group vertical-equity diagnostic defined above.}}}

{\color{impGreen}

\todo{Generate the two panels in Figure~\ref{fig:baseline_value_groups} from
the same held-out predictions used in Table~\ref{tab:ccao_baseline_results}.
Within each panel, construct the value groups using that model's
$V_i^{\mathrm{proxy}}$ from Eq.~\eqref{eq:market_value_proxy}. Report the
median ratio and its 90\% confidence interval for each group. Use identical
axes across panels and show the model's overall median ratio as the horizontal
reference.}


\begin{figure}[!ht]
    \centering
    \safeincludegraphics[width=0.8\textwidth]
    {paper/img/baseline_models_motivation_2024_2025_v3.pdf}
    \caption{
    Assessment ratios over sale price for the two baseline models in the held-out test and
    2025. The horizontal axis uses a base-10 logarithmic scale. Curves show
    median ratios in 30 equal-count sale-price bins, shaded regions show the
    corresponding interquartile ranges, and light points are a fixed random
    sample of 2,500 sales per panel. All panels use identical axes and the
    ratio-one reference line.
    }
    \label{fig:baseline_motivation}
\end{figure}

\begin{figure}[!ht]
    \centering
    \safeincludegraphics[width=0.8\textwidth]
    {paper/img/baseline_models_decile_ratios_2024_2025.pdf}
    \caption{
    Median assessment ratios by ordered property-value group for the two
    baseline models in the held-out test and 2025. Groups are model-specific deciles of the
    market-value proxy in Eq.~\eqref{eq:market_value_proxy}; intervals are
    90\% percentile bootstrap confidence intervals from 1,000 within-decile
    resamples. Dashed lines show each panel's overall median ratio and dotted
    lines show the ratio-one reference. All panels use identical axes.
    }
    \label{fig:baseline_value_groups}
\end{figure}



\begin{figure}[!ht]
    \centering

    \begin{subfigure}[t]{0.485\textwidth}
        \centering
        \safeincludegraphics[width=\linewidth]
        {paper/img/baseline_linear_value_groups.pdf}
        \caption{Linear regression.}
        \label{fig:baseline_value_groups_linear}
    \end{subfigure}
    \hfill
    \begin{subfigure}[t]{0.485\textwidth}
        \centering
        \safeincludegraphics[width=\linewidth]
        {paper/img/baseline_lgbm_value_groups.pdf}
        \caption{Unpenalized LightGBM.}
        \label{fig:baseline_value_groups_lgbm}
    \end{subfigure}

    \caption{
    Assessment-ratio profiles across ordered property-value groups for the
    two baseline models. Points show group median ratios and intervals show
    the corresponding 90\% confidence intervals. Within each panel, groups
    are ordered from lower to higher values using that model's market-value
    proxy in Eq.~\eqref{eq:market_value_proxy}; the horizontal reference is
    the model's overall median ratio. Identical axes are used across panels.
    }
    \label{fig:baseline_value_groups}
\end{figure}

}


% -------------------------------------------------------------------------
% Interpretation of the baseline evidence
% -------------------------------------------------------------------------

{\color{impGreen}

Table~\ref{tab:ccao_baseline_results} and
Figure~\ref{fig:baseline_value_groups} apply the evaluation framework in
Section~\ref{subsec:metrics} to the two baseline model classes. Assessment
level and vertical equity are interpreted separately: the former concerns
aggregate calibration of the ratios, whereas the latter concerns how those
ratios vary across property-value levels.

\todo{After the final empirical run, replace this note with one concise
results paragraph. Describe only the comparisons needed for the motivation:
(i) the change in predictive performance, (ii) the change in aggregate ratio
dispersion, (iii) assessment level as a separate dimension, and (iv) whether
the complementary vertical-equity measures and the value-group profiles show
a stronger regressive pattern under either baseline. Do not narrate every
entry in Table~\ref{tab:ccao_baseline_results}.}

}

% -------------------------------------------------------------------------
% Scope of the comparison and empirical motivation
% -------------------------------------------------------------------------

\textcolor{blue}{\sout{The baseline comparison should not be interpreted as evidence that machine-learning models are intrinsically more regressive than linear models.}}

\textcolor{blue}{\sout{Instead, the purpose of the comparison is to identify the application-specific tension motivating the proposed method: improvements in predictive accuracy and aggregate ratio dispersion need not eliminate systematic price-related regressivity. The assessor-facing vertical-equity suite and supplemental robustness diagnostics defined above remain the out-of-sample criteria for evaluating that behavior.}}

{\color{impGreen}

The comparison is application-specific and should not be interpreted as
evidence that machine-learning or nonlinear valuation models are intrinsically
more regressive than linear models. Its purpose is to determine whether the
predictive advantages of the more flexible baseline are accompanied by
stronger departures from vertical neutrality under the common evaluation
framework.

\todo{Once the final results are populated, make the next statement
declarative and consistent with the evidence in
Table~\ref{tab:ccao_baseline_results} and
Figure~\ref{fig:baseline_value_groups}.}

If the more flexible baseline provides meaningful predictive or ratio-dispersion
gains while exhibiting a stronger value-related vertical-equity pattern, the
model-design question is not simply whether to return to the linear benchmark.
It is whether those predictive advantages can be retained while reducing the
value-related pattern through the training objective.

}


% -------------------------------------------------------------------------
% Bridge to the proposed method
% -------------------------------------------------------------------------

\textcolor{blue}{\sout{Rather than equating any one of those metrics with the training objective, the method uses the identity $e_i(f)=\log r_i(f)$ together with $y_i=\log P_i$ to define a tractable empirical first-order target.}}

\textcolor{blue}{\sout{$\widehat{\operatorname{Cov}}\bigl(e(f),y\bigr)
=
\widehat{\operatorname{Cov}}\bigl(\log r(f),\log P\bigr).$}}

\textcolor{blue}{\sout{This covariance is a trainable first-order proxy for value-related residual structure, not an assessor-standard definition of vertical equity. A covariance near zero can coexist with nonlinear, distributional, or localized value-related patterns. Section~\ref{sec:method} therefore uses covariance as the training mechanism while the broader assessor-facing and robustness suite above provides the independent validation of whether that mechanism actually reduces regressivity out of sample.}}

{\color{impGreen}

The assessor-facing diagnostics above define how the resulting valuations are
evaluated; they need not be optimized directly during training. By
Eqs.~\eqref{eq:prediction_scale} and~\eqref{eq:log_ratio},
\[
e_i(f)=\log r_i(f),
\qquad
y_i=\log P_i.
\]
A first-order value-related pattern can therefore be represented on the
training scale through
\begin{equation}
\label{eq:training_covariance_bridge}
\widehat{\operatorname{Cov}}\bigl(e(f),y\bigr)
=
\widehat{\operatorname{Cov}}\bigl(\log r(f),\log P\bigr).
\end{equation}
Covariance is thus used as a tractable training proxy rather than as a
replacement for the assessor-facing vertical-equity measures.
Section~\ref{sec:method} develops this correction.

}



\section{Covariance-Guided Regressivity Correction}
\label{sec:method}

% \textcolor{blue}{\sout{This section present translates the vertical-equity concern identified in Section~\ref{sec:setting} into a trainable objective. The construction proceeds in five steps. We first express assessment ratios as residuals on the log-price scale and use covariance to measure their global linear relationship with property values. We then introduce a direct squared-covariance penalty and derive a sample-additive, covariance-guided surrogate. Next, we instantiate both formulations for linear regression and second-order tree boosting. Finally, we clarify what the penalties can and cannot control and describe how the penalty level is selected using temporal validation and assessor-facing diagnostics.}}

{This section translates the value-related regressivity identified in Section~\ref{sec:setting} into a training-time correction penalty. We first define a global first-order covariance-based regressivity target on the log-price scale, then introduce a direct squared-covariance penalty, and finally derive a sample-additive upper-bound surrogate that is compatible with standard boosting software. The correction penalty leaves the prediction target, feature information, and base model class unchanged. That is, for a fixed penalty formulation, the only additional tuning parameter is the penalty strength, which we denote $\rho$.}

{Let $J_0(f)$ denote the baseline training objective, whose data-fit component is the squared log-price loss $L(f)$ defined in Section~\ref{subsec:baseline_models}. All quantities below are computed from the ``current'' training sample. In temporal validation, the centering quantities and penalties are recomputed using the training block of each fold only.}


\subsection{{Log-Scale Regressivity Target}}
\label{subsec:training_regressivity}

Let
\begin{equation}
\label{eq:centered_log_price}
\bar y=\frac{1}{n}\sum_{i=1}^n y_i,
\qquad
c_i=y_i-\bar y,
\qquad
\bar e(f)=\frac{1}{n}\sum_{i=1}^n e_i(f).
\end{equation}
The empirical covariance between log-residuals and log-prices is
\begin{equation}
\label{eq:training_covariance}
C(f)
=
\widehat{\operatorname{Cov}}\bigl(e(f),y\bigr)
=
\frac{1}{n}\sum_{i=1}^n
\bigl(e_i(f)-\bar e(f)\bigr)(y_i-\bar y)
=
\frac{1}{n}\sum_{i=1}^n e_i(f)c_i.
\end{equation}
The last equality follows from $\sum_{i=1}^n c_i=0$, so centering the residuals does not change the covariance.

This covariance has a direct slope interpretation. Let
\[
V_y=\frac{1}{n}\sum_{i=1}^n c_i^2.
\]
Provided $V_y>0$, the least-squares slope from regressing log assessment ratios on log prices is
\begin{equation}
\label{eq:covariance_slope}
\operatorname{slope}\bigl(e(f)\sim y\bigr)
=
\frac{C(f)}{V_y}.
\end{equation}
Because $V_y$ is fixed by the training sample, reducing $|C(f)|$ is equivalent to reducing the magnitude of this global linear trend. In particular, $C(f)<0$ indicates that log assessment ratios tend to decline with log sale price, which is the first-order pattern associated with regressivity.

{This slope is a training-scale mechanism measure, not the assessor-facing PRB statistic defined in Section~\ref{subsec:metrics}. We use covariance rather than correlation because $V_y$ is fixed within the training sample, whereas correlation additionally normalizes by the model-dependent residual variance. Covariance therefore provides a direct training target for the global value-related slope.}

{Sale price enters this construction as the observed market-value benchmark for verified training sales; it is not introduced as a protected attribute. The broader connection to dependence-based fair regression is discussed in Section~\ref{subsec:related_work}.}


\subsection{{Direct Squared-Covariance Penalty}}
\label{subsec:direct_penalty}

% \textcolor{blue}{\sout{We therefore define the direct covariance-penalized learner as}}
{We first penalize the chosen first-order target directly. The \emph{covariance-penalized learner} is}

% \begin{center}
% % \textcolor{blue}{\sout{\mbox{$\displaystyle
% % f_\rho^{\mathrm{cov}}
% % \in
% % \argminop_{f\in\Hcal}
% \left\{
% L(f)+\Omega(f)+\frac{\rho}{2}C(f)^2
% \right\},
% \qquad
% \rho\ge 0.
% $}}}
% \end{center}

{
\begin{equation}
\label{eq:direct_penalty}
f_\rho^{\mathrm{cov}}
\in
\argminop_{f\in\Hcal}
\left\{
J_0(f)+\frac{\rho}{2}C(f)^2
\right\},
\qquad
\rho\geq 0.
\end{equation}
}

When $\rho=0$, formulation~\eqref{eq:direct_penalty} reduces to the baseline learner. As $\rho$ increases, the optimization places greater weight on reducing the global residual--price trend. Squaring the covariance treats negative and positive trends symmetrically. This is important because the objective is not to replace regressivity with progressivity, but to move the first-order relationship toward zero.

% \textcolor{blue}{\sout{The direction of the correction can be seen from the gradient. Let $e=(e_1,\ldots,e_n)^\top$ and $c=(c_1,\ldots,c_n)^\top$. Excluding $\Omega(f)$, the gradient and Hessian with respect to the prediction vector $\widehat y$ are}}

% \begin{center}
% \textcolor{blue}{\sout{\mbox{$\displaystyle
% \nabla_{\widehat y}
% \left(L(f)+\frac{\rho}{2}C(f)^2\right)
% =
% \frac{2}{n}e+\frac{\rho}{n}C(f)c
% $}}}

% \textcolor{blue}{\sout{\mbox{$\displaystyle
% \nabla_{\widehat y}^2
% \left(L(f)+\frac{\rho}{2}C(f)^2\right)
% =
% \frac{2}{n}I+\frac{\rho}{n^2}cc^\top.
% $}}}
% \end{center}

{For interpretation, it is enough to consider the contribution of the covariance penalty to the gradient of observation $i$:}
{
\begin{equation}
\label{eq:direct_penalty_gradient}
\frac{\partial}{\partial \widehat y_i}
\left[
\frac{\rho}{2}C(f)^2
\right]
=
\frac{\rho}{n}C(f)c_i.
\end{equation}
}

% \textcolor{blue}{\sout{Suppose, for example, that the fitted model is regressive, so $C(f)<0$. For a lower-value property, $c_i<0$, and the covariance component of the gradient is positive; a descent step therefore lowers its predicted value. For a higher-value property, $c_i>0$, and the covariance component is negative; a descent step raises its predicted value. The penalty consequently acts in the direction needed to flatten the global ratio--value relationship. The direction reverses automatically when the model is progressive.}}


% \textcolor{blue}{\sout{Suppose, for example, that the fitted model is regressive due to $C(f)<0$. For a sale below the geometric-mean price (lower-tier property), $c_i<0$, and thus the covariance component of the gradient is positive, i.e., a descent step lowers its predicted value. For a sale above the geometric-mean price (higher-tier property), $c_i>0$, and thus the covariance component is negative, i.e., a descent step raises its predicted value. The penalty therefore acts in the direction required to flatten the global log-ratio--log-price trend. Note that the direction and interpretation reverse automatically when the fitted trend is progressive.}}

{Suppose, for example, that the fitted model is regressive due to $C(f)<0$. For a sale below the geometric-mean price, $c_i<0$, and thus the covariance component of the gradient is positive, i.e., a descent step lowers its predicted value. For a sale above the geometric-mean price, $c_i>0$, and thus the covariance component is negative, i.e., a descent step raises its predicted value. The penalty therefore acts in the direction required to flatten the global log-ratio--log-price trend. Note that the direction and interpretation reverse automatically when the fitted trend is progressive.}


Note that the direct penalty targets how assessment ratios vary with sale price, but not their overall assessment level. Multiplying every predicted value by the same positive constant shifts all log assessment ratios by the same amount but leaves the covariance $C(f)$ unchanged. Thus, even a model with $C(f)=0$ can have a median or mean assessment ratio different from one. Hence, assessment level must be evaluated separately.

% {The direct penalty controls this first-order slope separately from overall assessment level. Adding the same constant to every log prediction leaves $C(f)$ unchanged, i.e., multiplying every price prediction by the same positive constant does not change the targeted covariance. Assessment-level calibration therefore remains a separate modeling and evaluation dimension.}

% \textcolor{blue}{\sout{The direct formulation most faithfully represents the statistical target, but it creates a computational complication. The rank-one term $cc^\top$ in~\eqref{eq:direct_gradient_hessian} couples all observations. Standard second-order boosting interfaces instead expect one gradient and one scalar Hessian for each observation. Section~\ref{subsec:model_specific_implementations} therefore develops both a diagonal Hessian approximation for standard boosting software and an exact leaf-weight solver for a fixed tree structure.}}

{For second-order models (e.g., LightGBM), the corresponding Hessian of the squared log-price loss plus direct penalty is}
{
\begin{equation}
\label{eq:direct_gradient_hessian}
\nabla_{\widehat y}^2
\left(
L(f)+\frac{\rho}{2}C(f)^2
\right)
=
\frac{2}{n}I+\frac{\rho}{n^2}cc^\top.
\end{equation}
}
{The rank-one Hessian term $\frac{\rho}{n^{2}}cc^\top$ introduces cross-observation curvature. Standard LightGBM custom objectives provide one gradient and one Hessian value per observation, so this dense term cannot be represented exactly through the usual interface. An exact second-order implementation therefore requires additional machinery: using only the diagonal is possible, but discards the cross-observation curvature. The next subsection proposes an alternative penalty term to handle this caveat. Further, the model-specific derivations are given in Appendix~\ref{app:implementation}.


% . The exact second-order implementation therefore requires additional machinery: a diagonal approximation is possible but discards these cross-observation terms. The model-specific derivations are given in Appendix~\ref{app:implementation}.}

% \textcolor{blue}{\sout{The fixed-prediction-space analysis in Appendix~\ref{app:theory} provides additional intuition. Within a fixed linear prediction space, the magnitude of the covariance decreases smoothly with $\rho$, while the increase in in-sample squared error is second order near $\rho=0$. This explains why a modest penalty can, in principle, meaningfully reduce the targeted trend at a relatively small accuracy cost. For boosted trees, this result is only a local benchmark because changing $\rho$ can also change splits, leaf structures, and stopping behavior.}}


\subsection{{Sample-Additive Covariance Upper-Bound Surrogate}}
\label{subsec:surrogate}

% \textcolor{blue}{\sout{The direct penalty presents two practical limitations. First, covariance is an aggregate statistic, so positive and negative local relationships can offset one another. Second, its squared form is not additive across observations and therefore does not fit directly into the standard custom-objective interface used by LightGBM. We address the computational limitation by deriving a sample-additive surrogate motivated by an upper bound on the squared covariance.}}

{For standard second-order boosting software, we seek a penalty that decomposes exactly across observations. Given that~\eqref{eq:training_covariance} already gives $C(f)=\frac{1}{n}\sum_{i=1}^n e_i(f)c_i$, we can make use of Jensen's inequality that yields a sample-additive upper bound to the squared-covariance term directly.}

% \textcolor{blue}{\sout{By Jensen's inequality,}}
% \begin{center}
% \textcolor{blue}{\sout{\mbox{$\displaystyle
% C(f)^2
% =
% \left[
% \frac{1}{n}\sum_{i=1}^n
% \bigl(e_i(f)-\bar e(f)\bigr)c_i
% \right]^2
% \leq
% \frac{1}{n}\sum_{i=1}^n
% \bigl(e_i(f)-\bar e(f)\bigr)^2c_i^2.
% $}}}
% \end{center}

{
\begin{equation}
\label{eq:surrogate_upper_bound}
C(f)^2
=
\left(
\frac{1}{n}\sum_{i=1}^n e_i(f)c_i
\right)^2
\leq
\frac{1}{n}\sum_{i=1}^n e_i(f)^2c_i^2.
\end{equation}
}
% \textcolor{blue}{\sout{Although the right-hand side is a sum, it remains coupled through the model-dependent mean residual $\bar e(f)$. To obtain a fully observation-additive objective, we replace the centered residual $e_i(f)-\bar e(f)$ with $e_i(f)$ and define}}
{This inequality motivates the fully observation-additive surrogate penalty}
\begin{equation}
\label{eq:surrogate}
\Psi^{\mathrm{surr}}(f)
=
\frac{1}{n}\sum_{i=1}^n e_i(f)^2c_i^2
=
\frac{1}{n}\sum_{i=1}^n
e_i(f)^2(y_i-\bar y)^2.
\end{equation}
The resulting \emph{surrogate-penalized learner} is therefore
\begin{equation}
\label{eq:surrogate_learning}
f_\rho^{\mathrm{surr}}
\in
\argminop_{f\in\Hcal}
\left\{
{J_0(f)}
+
\rho\Psi^{\mathrm{surr}}(f)
\right\}.
\end{equation}
% \textcolor{blue}{\sout{After replacing the centered residual, $\Psi^{\mathrm{surr}}$ is not, in general, an upper bound on $C(f)^2$. It should therefore be interpreted as a \emph{covariance-guided} regularizer rather than exact covariance control. Its connection to the direct penalty is that both emphasize residuals according to their position in the value distribution. The surrogate makes this principle observation-specific and computationally separable.}}
% \textcolor{blue}{\sout{Thus, $\Psi^{\mathrm{surr}}(f)$ is a genuine upper bound on the squared empirical covariance in~\eqref{eq:training_covariance}. Nevertheless, the bound does not make surrogate minimization equivalent to direct covariance minimization: the upper bound can be loose when the signed contributions $e_i(f)c_i$ cancel in aggregate covariance. The surrogate instead penalizes the squared magnitude of each contribution before aggregation, so opposite-signed contributions cannot cancel within the penalty.}}
{
Thus, $\Psi^{\mathrm{surr}}(f)$ is a genuine upper bound on the squared empirical covariance in~\eqref{eq:training_covariance}. Its exact relationship with the direct-covariance penalty can be seen by making the observation
% \[
% q_i(f)=e_i(f)c_i.
% \]
% Since $C(f)=n^{-1}\sum_i q_i(f)$,
\begin{equation}
\label{eq:surrogate_decomposition}
\Psi^{\mathrm{surr}}(f)
=
C(f)^2
+
\frac{1}{n}
\sum_{i=1}^n
\left(e_i(f)c_i-C(f)\right)^2.
\end{equation}
Hence, the surrogate controls squared covariance together with the dispersion of the observation-level residual--price contributions. It is therefore generally more conservative than direct covariance minimization: opposite-signed contributions can cancel in $C(f)$ but not in $\Psi^{\mathrm{surr}}(f)$.
}


{Thus, $\Psi^{\mathrm{surr}}(f)$ is a genuine upper bound on the squared empirical covariance in~\eqref{eq:training_covariance}. Nevertheless, the bound does not make surrogate minimization equivalent to direct covariance minimization: the upper bound can be loose when the signed contributions $e_i(f)c_i$ cancel in aggregate covariance. The surrogate instead penalizes the squared magnitude of each contribution before aggregation, so opposite-signed contributions cannot cancel within the penalty.}

Combining the baseline loss with the surrogate gives
\begin{equation}
\label{eq:surrogate_weighted_loss}
L(f)+\rho\Psi^{\mathrm{surr}}(f)
=
\frac{1}{n}\sum_{i=1}^n
\left[1+\rho(y_i-\bar y)^2\right]e_i(f)^2.
\end{equation}
Thus, the surrogate is equivalent to a weighted squared-error objective with fixed observation weights
\begin{equation}
\label{eq:surrogate_weights}
w_i(\rho)=1+\rho(y_i-\bar y)^2.
\end{equation}
Properties near the center of the log-price distribution receive weights close to the baseline weight, whereas properties in the lower and upper tails receive larger weights.

% \textcolor{blue}{\sout{Because $\bar y$ is the logarithm of the geometric mean sale price, the weighting is symmetric in proportional distance from the center of the price distribution. For example, equally sized multiplicative price deviations above and below this center receive the same additional weight.}}

% {Since $\exp(\bar y)$ is the geometric mean sale price, the weight function is symmetric in log-distance from this center: prices that differ from the geometric mean by reciprocal multiplicative factors receive the same additional weight. This symmetry concerns the weighting function only: the empirical lower and upper tails, and the resulting prediction changes, need not be symmetric.}

% \textcolor{blue}{\sout{Since $\exp(\bar y)$ is the geometric mean sale price, the weight function is symmetric in log-distance from this center: prices that differ from the geometric mean by reciprocal multiplicative factors receive the same additional weight. This symmetry concerns the weighting function only: the empirical lower and upper tails, and the resulting prediction changes, need not be symmetric.}}

{Since $\exp(\bar y)$ is the geometric mean sale price, the weight function is symmetric in log-distance from this center: prices that differ from the geometric mean by reciprocal multiplicative factors receive the same additional weight. This symmetry concerns the weighting function only: the empirical lower and upper tails, and the resulting prediction changes, need not be symmetric. Moreover, because the additional weight grows quadratically with log-price distance, observations far from the center can receive disproportionate influence. The log transformation compresses this leverage relative to an analogous raw-price construction, but does not bound it.}

% \textcolor{blue}{\sout{This representation clarifies the mechanism without imposing a directional assumption on individual observations. The surrogate raises the cost of residuals farther from the center of the log-price distribution. Because it uses $e_i(f)^2c_i^2$, however, it does not encode whether the aggregate residual--price trend is regressive or progressive. Its effect on covariance and on assessor-facing vertical-equity diagnostics must therefore be evaluated out of sample.}}

{This representation clarifies the mechanism without imposing a directional assumption on individual observations. The surrogate raises the cost of residuals farther from the center of the log-price distribution. Unlike the direct penalty, its observation-level correction is not determined by the sign of the aggregate covariance: it depends on the residual $e_i(f)$ and the squared log-price deviation $c_i^2$. Its resulting effect on the signed covariance and on assessor-facing vertical-equity diagnostics must therefore be evaluated out of sample.}

% \textcolor{blue}{\sout{This representation also clarifies why the surrogate can affect the ratio--value trend. Under a regressive baseline, the largest systematic residuals often occur toward the lower and upper ends of the value distribution. By increasing the cost of errors in both tails, the surrogate discourages the model from obtaining strong average predictive performance by allowing systematically positive residuals at the low end and negative residuals at the high end. Unlike covariance itself, the pointwise squared terms cannot cancel across observations.}}

% {This representation clarifies the mechanism without imposing a directional assumption on individual observations. The surrogate raises the cost of residuals farther from the center of the log-price distribution. Because it uses $e_i(f)^2c_i^2$, however, it does not encode whether the aggregate residual--price trend is regressive or progressive. Its effect on covariance and on assessor-facing vertical-equity diagnostics must therefore be evaluated out of sample.}

For squared loss, the observation-wise gradient and Hessian are
\begin{align}
\label{eq:surrogate_derivatives}
g_i
&=
\frac{2}{n}e_i(f)
\left[1+\rho(y_i-\bar y)^2\right],
\\
h_i
&=
\frac{2}{n}
\left[1+\rho(y_i-\bar y)^2\right].
\end{align}
Both quantities depend only on observation $i$ once $\bar y$ has been computed. They can therefore be supplied through a standard second-order boosting custom-objective interface.

% \textcolor{blue}{\sout{The implementation uses the corresponding convention with the common normalization absorbed into the objective scaling.}}

{The implementation follows LightGBM's observation-wise derivative convention. The exact numerical scaling of $\rho$ is therefore tied to the implemented objective parameterization and is documented with the empirical specification.}

% \textcolor{blue}{\sout{The direct and surrogate penalties consequently serve different methodological roles. The direct penalty is the clearest representation of the statistical quantity being controlled, but its curvature couples observations. The surrogate sacrifices exact covariance control in exchange for sample additivity, standard-software compatibility, and reduced sensitivity to cancellation across different parts of the value distribution. Moreover, the numerical values of $\rho$ under the two formulations are not directly comparable because the penalties have different scales. In the CCAO application, the surrogate provides the strongest combined validation performance and is therefore the principal practical specification.}}

{Unlike the direct covariance penalty, the surrogate is not invariant to a common shift in all log predictions because it penalizes the magnitude of $e_i(f)$ itself. It can therefore affect overall assessment level in addition to the value-related trend. %Assessment-level diagnostics consequently remain part of validation and out-of-sample evaluation.
}


\subsection{\textcolor{impGreen}{Interpretation, Validation, and Scope}}
\label{subsec:correction_scope}

% \textcolor{blue}{\sout{The two formulations serve complementary roles. The direct penalty targets the chosen global first-order association between log assessment ratios and log sale prices. The surrogate replaces direct covariance minimization with a sample-additive upper-bound objective that is directly compatible with standard boosting software. Because the two penalties have different scales, the numerical values of $\rho$ are not directly comparable across formulations. Their appropriate scale can also vary across training samples and jurisdictions, so $\rho$ is calibrated separately through chronological validation.}}

\textcolor{impGreen}{The two formulations intervene differently. The direct penalty responds to the aggregate signed first-order association between log assessment ratios and log sale prices. The surrogate instead reallocates squared-error importance across observations according to their distance from the center of the log-price distribution. Although this weighted objective upper-bounds squared covariance, it generally controls more than the targeted first-order trend, as shown in~\eqref{eq:surrogate_decomposition}. Because the two penalties have different scales, the numerical values of $\rho$ are not directly comparable across formulations. Their appropriate scale can also vary across training samples and jurisdictions, so $\rho$ is calibrated separately through chronological validation.}

% \textcolor{impGreen}{The two formulations serve complementary roles. The direct penalty targets the chosen global first-order association between log assessment ratios and log sale prices. The surrogate replaces direct covariance minimization with a sample-additive upper-bound objective that is directly compatible with standard boosting software. Because the two penalties have different scales, the numerical values of $\rho$ are not directly comparable across formulations. Their appropriate scale can also vary across training samples and jurisdictions, so $\rho$ is calibrated separately through chronological validation.}

\textcolor{impGreen}{Neither formulation directly optimizes PRD, PRB, VEI, COD, or assessment level. Nor does zero training covariance imply the absence of nonlinear, geographic, property-class, temporal, or other local value-related patterns. The training covariance is therefore treated as a mechanism target rather than as a sufficient measure of vertical equity. Predictive performance, assessment level, PRD, PRB, VEI, ratio curves, and subgroup diagnostics remain separate validation and out-of-sample evaluation criteria.}

\textcolor{impGreen}{Sale price is required to construct the correction only for verified sales used during model fitting. Once a model has been trained, the resulting valuation function continues to generate predictions from $x$ alone, as in the baseline AVM; the correction does not require a current sale price for an unsold property.}

\textcolor{impGreen}{Section~\ref{sec:empirical_design} specifies the chronological validation and penalty-selection procedure used in the empirical application. Appendix~\ref{app:theory} provides complementary intuition for the direct penalty: within a fixed linear prediction space, increasing $\rho$ shrinks the targeted covariance smoothly while the increase in squared error is second order near $\rho=0$. This result is an interpretive benchmark, not a guarantee for fully retrained boosted trees, whose splits, leaves, and stopping behavior can change with the penalty.}




% \section{Covariance-Guided Regressivity Correction}
% \label{sec:method}

% This section present translates the vertical-equity concern identified in Section~\ref{sec:setting} into a trainable objective. The construction proceeds in five steps. We first express assessment ratios as residuals on the log-price scale and use covariance to measure their global linear relationship with property values. We then introduce a direct squared-covariance penalty and derive a sample-additive, covariance-guided surrogate. Next, we instantiate both formulations for linear regression and second-order tree boosting. Finally, we clarify what the penalties can and cannot control and describe how the penalty level is selected using temporal validation and assessor-facing diagnostics.


% \subsection{Direct Squared-Covariance Penalty}
% \label{subsec:direct_penalty}

% Let
% \begin{equation}
% \label{eq:centered_log_price}
% \bar y=\frac{1}{n}\sum_{i=1}^n y_i,
% \qquad
% c_i=y_i-\bar y,
% \qquad
% \bar e(f)=\frac{1}{n}\sum_{i=1}^n e_i(f).
% \end{equation}
% The empirical covariance between log residuals and log prices is
% \begin{align}
% \label{eq:training_covariance}
% C(f)=
% \widehat{\Cov}\bigl(e(f),y\bigr) =
% \frac{1}{n}\sum_{i=1}^n
% \bigl(e_i(f)-\bar e(f)\bigr)(y_i-\bar y) =
% \frac{1}{n}\sum_{i=1}^n e_i(f)c_i.
% \end{align}
% The last equality follows from \(\sum_i c_i=0\), so centering the residuals does not change the covariance.

% This covariance has a direct slope interpretation. Let
% \[
% V_y=\frac{1}{n}\sum_{i=1}^n c_i^2.
% \]
% Provided \(V_y>0\), the least-squares slope from regressing log assessment ratios on log prices is
% \begin{equation}
% \label{eq:covariance_slope}
% \operatorname{slope}\bigl(e(f)\sim y\bigr)
% =
% \frac{C(f)}{V_y}.
% \end{equation}
% Because \(V_y\) is fixed by the training sample, reducing \(|C(f)|\) is equivalent to reducing the magnitude of this global linear trend. In particular, \(C(f)<0\) indicates that log assessment ratios tend to decline with log price, which is the first-order pattern associated with regressivity.

% We therefore define the direct covariance-penalized learner as
% \begin{equation}
% \label{eq:direct_penalty}
% f_\rho^{\mathrm{cov}}
% \in
% \argminop_{f\in\Hcal}
% \left\{
% L(f)+\Omega(f)+\frac{\rho}{2}C(f)^2
% \right\},
% \qquad
% \rho\ge 0.
% \end{equation}
% When \(\rho=0\), formulation~\eqref{eq:direct_penalty} reduces to the baseline learner. As \(\rho\) increases, the optimization places greater weight on reducing the global residual--price trend. Squaring the covariance treats negative and positive trends symmetrically. This is important because the objective is not to replace regressivity with progressivity, but to move the first-order relationship toward zero.

% The direction of the correction can be seen from the gradient. Let \(e=(e_1,\ldots,e_n)^\top\) and \(c=(c_1,\ldots,c_n)^\top\). Excluding \(\Omega(f)\), the gradient and Hessian with respect to the prediction vector \(\widehat y\) are
% \begin{align}
% \label{eq:direct_gradient_hessian}
% \nabla_{\widehat y}
% \left(
% L(f)+\frac{\rho}{2}C(f)^2
% \right)
% &=
% \frac{2}{n}e+\frac{\rho}{n}C(f)c,
% \\
% \nabla_{\widehat y}^2
% \left(
% L(f)+\frac{\rho}{2}C(f)^2
% \right)
% &=
% \frac{2}{n}I+\frac{\rho}{n^2}cc^\top.
% \end{align}
% Suppose, for example, that the fitted model is regressive, so \(C(f)<0\). For a lower-value property, \(c_i<0\), and the covariance component of the gradient is positive; a descent step therefore lowers its predicted value. For a higher-value property, \(c_i>0\), and the covariance component is negative; a descent step raises its predicted value. The penalty consequently acts in the direction needed to flatten the global ratio--value relationship. The direction reverses automatically when the model is progressive.

% The direct formulation most faithfully represents the statistical target, but it creates a computational complication. The rank-one term \(cc^\top\) in~\eqref{eq:direct_gradient_hessian} couples all observations. Standard second-order boosting interfaces instead expect one gradient and one scalar Hessian for each observation. Section~\ref{subsec:model_specific_implementations} therefore develops both a diagonal Hessian approximation for standard boosting software and an exact leaf-weight solver for a fixed tree structure.

% The fixed-prediction-space analysis in Appendix~\ref{app:theory} provides additional intuition. Within a fixed linear prediction space, the magnitude of the covariance decreases smoothly with \(\rho\), while the increase in in-sample squared error is second order near \(\rho=0\). This explains why a modest penalty can, in principle, meaningfully reduce the targeted trend at a relatively small accuracy cost. For boosted trees, this result is only a local benchmark because changing \(\rho\) can also change splits, leaf structures, and stopping behavior.

% \subsection{Sample-Additive Covariance-Guided Surrogate}
% \label{subsec:surrogate}

% The direct penalty presents two practical limitations. First, covariance is an aggregate statistic, so positive and negative local relationships can offset one another. Second, its squared form is not additive across observations and therefore does not fit directly into the standard custom-objective interface used by LightGBM. We address the computational limitation by deriving a sample-additive surrogate motivated by an upper bound on the squared covariance.

% By Jensen's inequality,
% \begin{align}
% \label{eq:surrogate_upper_bound}
% C(f)^2
% &=
% \left[
% \frac{1}{n}\sum_{i=1}^n
% \bigl(e_i(f)-\bar e(f)\bigr)c_i
% \right]^2 \le
% \frac{1}{n}\sum_{i=1}^n
% \bigl(e_i(f)-\bar e(f)\bigr)^2c_i^2.
% \end{align}
% Although the right-hand side is a sum, it remains coupled through the model-dependent mean residual \(\bar e(f)\). To obtain a fully observation-additive objective, we replace the centered residual \(e_i(f)-\bar e(f)\) with \(e_i(f)\) and define
% \begin{equation}
% \label{eq:surrogate}
% \Psi^{\mathrm{surr}}(f)
% =
% \frac{1}{n}\sum_{i=1}^n e_i(f)^2c_i^2
% =
% \frac{1}{n}\sum_{i=1}^n
% e_i(f)^2(y_i-\bar y)^2.
% \end{equation}
% The resulting learner is
% \begin{equation}
% \label{eq:surrogate_learning}
% f_\rho^{\mathrm{surr}}
% \in
% \argminop_{f\in\Hcal}
% \left\{
% L(f)+\Omega(f)+\rho\Psi^{\mathrm{surr}}(f)
% \right\}.
% \end{equation}

% After replacing the centered residual, \(\Psi^{\mathrm{surr}}\) is not, in general, an upper bound on \(C(f)^2\). It should therefore be interpreted as a \emph{covariance-guided} regularizer rather than exact covariance control. Its connection to the direct penalty is that both emphasize residuals according to their position in the value distribution. The surrogate makes this principle observation-specific and computationally separable.

% Combining the baseline loss with the surrogate gives
% \begin{equation}
% \label{eq:surrogate_weighted_loss}
% L(f)+\rho\Psi^{\mathrm{surr}}(f)
% =
% \frac{1}{n}\sum_{i=1}^n
% \left[1+\rho(y_i-\bar y)^2\right]e_i(f)^2.
% \end{equation}
% Thus, the surrogate is equivalent to a weighted squared-error objective with fixed observation weights
% \[
% w_i(\rho)=1+\rho(y_i-\bar y)^2.
% \]
% Properties near the center of the log-price distribution receive weights close to the baseline weight, whereas properties in the lower and upper tails receive larger weights. Because \(\bar y\) is the logarithm of the geometric mean sale price, the weighting is symmetric in proportional distance from the center of the price distribution. For example, equally sized multiplicative price deviations above and below this center receive the same additional weight.

% This representation also clarifies why the surrogate can affect the ratio--value trend. Under a regressive baseline, the largest systematic residuals often occur toward the lower and upper ends of the value distribution. By increasing the cost of errors in both tails, the surrogate discourages the model from obtaining strong average predictive performance by allowing systematically positive residuals at the low end and negative residuals at the high end. Unlike covariance itself, the pointwise squared terms cannot cancel across observations.

% For squared loss, the observation-wise gradient and Hessian are
% \begin{align}
% \label{eq:surrogate_derivatives}
% g_i
% &=
% \frac{2}{n}e_i(f)
% \left[1+\rho(y_i-\bar y)^2\right],
% \\
% h_i
% &=
% \frac{2}{n}
% \left[1+\rho(y_i-\bar y)^2\right].
% \end{align}
% Both quantities depend only on observation \(i\) once \(\bar y\) has been computed. They can therefore be supplied directly through a standard second-order boosting custom-objective interface. The implementation uses the corresponding convention with the common normalization absorbed into the objective scaling.

% The direct and surrogate penalties consequently serve different methodological roles. The direct penalty is the clearest representation of the statistical quantity being controlled, but its curvature couples observations. The surrogate sacrifices exact covariance control in exchange for sample additivity, standard-software compatibility, and reduced sensitivity to cancellation across different parts of the value distribution. Moreover, the numerical values of \(\rho\) under the two formulations are not directly comparable because the penalties have different scales. In the CCAO application, the surrogate provides the strongest combined validation performance and is therefore the principal practical specification.












\section{Application and Evaluation Design}
\label{sec:empirical_design}

{This section specifies the empirical design used to evaluate the proposed correction. We first describe the CCAO application, data, and implementation; then define the chronological validation and model-selection protocol; next summarize the model comparators; and finally describe the exploratory six-county benchmark. All model comparisons reported in the next section use the evaluation criteria defined in Section~\ref{subsec:metrics}.}


\subsection{{The CCAO Residential Pipeline: Application, Data, and Implementation}}
\label{subsec:ccao_design}

% \textcolor{blue}{\sout{For the empirical analysis, we implement in Python the components of CCAO's residential valuation workflow required for model development and evaluation. The implementation closely follows CCAO's data-preparation, model-fitting, and ratio-study procedures, preserving the operational structure of the existing workflow while enabling controlled evaluation of alternative training objectives.}}

{For the primary application, we implement in Python the components of CCAO's residential valuation workflow required for model development and evaluation. The implementation follows the data, feature, model-fitting, and ratio-study structure of CCAO's published residential pipeline, but it is a research translation rather than a direct production run. The reported results should therefore be interpreted as evidence within this translated workflow, not as official CCAO assessment outputs.}

% \textcolor{blue}{\sout{The proposed methodology modifies only the model-training objective. We hold fixed the sale filters, predictor set, temporal data splits, transformation of predictions to the price scale, ratio-study procedures, and surrounding review process. This restricted intervention serves two purposes. First, it isolates the effect of incorporating vertical equity into model estimation. Second, it demonstrates how the methodology can be integrated into an existing mass-appraisal workflow without redesigning the broader valuation system.}}

{Within the LightGBM comparison, the intervention is deliberately restricted to the training objective. We hold fixed the sale filters, log-price target, underlying predictor information, temporal design, transformation back to the price scale, and evaluation procedures. The selected LightGBM hyperparameter configuration is also held fixed across the standard and penalized fits, so $\rho$ is the only additional tuning dimension introduced by the correction. This design isolates the practical effect of changing the objective while preserving the surrounding valuation workflow.}

The CCAO application uses verified Cook County residential sales {from the CCAO residential modeling data used for the final evaluation. \todo{Replace this phrase with the exact final data extract/version and citation used for the 2024 and 2025 evaluations.}}
{The residential scope includes the property types represented in the CCAO residential pipeline, including single-family homes, townhomes, small multifamily properties, and condominiums.}
Following the office pipeline, we exclude multicard PINs and sales marked as outliers.

% \textcolor{blue}{\sout{The pre-2024 modeling sample contains 349,661 observations. Sales from 2016--2022 form the development period; 2023 is held out for final testing. The 2024 sales block is reserved for later assessment-year analysis and is not used for the results reported here.}}

{The sales universe is 2016--2024. The current training file contains 2016 sales, so the first observed sale is on 2016-01-01. The oldest 90\% of eligible 2016--2024 sales (344,610 observations, through 2023-11-09) form the development/CV pool used for model fitting and all model and penalty selection. The newest 10\% (38,290 observations, 2023-11-09 through 2024-12-31) form the primary held-out test. The 2025 block contains 26,641 sales and is reported separately as a secondary forward evaluation after the selected configuration is refit on all 382,900 eligible 2016--2024 sales; neither the test nor 2025 enters CV or model selection. Because earlier exploratory versions of the project examined 2024 outputs, we describe the test as a held-out evaluation rather than as an untouched or preregistered confirmatory test.}

\begin{table}[!ht]
\centering
% \color{impGreen}
\caption{Final CCAO temporal samples.}
\label{tab:ccao_samples}
\begin{tabular}{llr}
\toprule
Sample & Period & Observations \\
\midrule
Development / validation & 2016-01-01--2023-11-09 & 344,610 \\
Primary held-out test & 2023-11-09--2024-12-31 & 38,290 \\
Production training & 2016-01-01--2024-12-31 & 382,900 \\
Secondary forward evaluation & 2025-01-01--2025-12-29 & 26,641 \\
\bottomrule
\end{tabular}
\end{table}
{\footnotesize The current training file contains 2016 sales; observed sales begin on 2016-01-01.}

The model uses 95 predictors, including property characteristics, location and proximity variables, Census measures, time variables, parcel-shape measures, neighborhood identifiers, and other assessor variables. Twenty-three are treated as categorical.
% \textcolor{blue}{\sout{All predictors are constructed upstream; the penalized models use the same feature set as the LightGBM baseline.}}
{All predictors are constructed upstream by the CCAO pipeline. Standard and penalized LightGBM fits use the same 95 predictors, with the same 23 variables treated as categorical; the penalty introduces no additional predictors or feature engineering. The linear benchmark uses the same underlying predictor information with learner-specific representation where required.}

{Table~\ref{tab:feature_groups} summarizes the predictor information used by the models. The full feature list is fixed before model fitting and should be archived with the replication materials.}

{
% \color{impGreen}
\begin{table}[!ht]
\centering
\caption{CCAO predictor groups used in the empirical models.}
\label{tab:feature_groups}
\begin{tabular}{lrl}
\toprule
Feature group & Count & Examples \\
\midrule
Assessor/property variables & 30 & Class, building area, land area, age, rooms, neighborhood \\
Location variables & 9 & Coordinates, census tract, school district, municipality \\
Proximity variables & 26 & Transit, parks, roads, hospitals, water, airports \\
ACS 5-year measures & 16 & Income, education, poverty, employment, tenure \\
Time variables & 8 & Sale year, month, quarter, day, post-COVID indicator \\
Parcel-shape variables & 6 & Edge/angle dispersion, bounding ratios, vertex count \\
\midrule
Total & 95 & 23 categorical predictors in LightGBM \\
\bottomrule
\end{tabular}
\end{table}
}


\subsection{\textcolor{impGreen}{Temporal Validation and Model Selection}}
\label{subsec:selection}
\todo{MOVE TO THE APPENDIX}

% \textcolor{blue}{\sout{Within 2016--2022, we use eight expanding-window rolling-origin folds. Each validation block follows its training block in time. For the principal CCAO comparison, candidates are screened using the average validation PRD reference band and then ranked by predictive performance. PRB, VEI, assessment level, and dispersion are reported as additional diagnostics rather than used to select the displayed model. The 2023 test set is evaluated only after selection.}}

\textcolor{impGreen}{Within the 344,610-sale development/CV pool, we use seven expanding-window rolling-origin folds. The cumulative windows advance in 15-month steps anchored at 2016-01-01, and within each cumulative window the newest 10\% of observations form validation. For fold $v$, let $\mathcal T_v$ and $\mathcal V_v$ denote its training and validation samples. Every fold respects the chronological ordering}
{\color{impGreen}
\begin{equation}
\label{eq:temporal_ordering}
\max_{i\in\mathcal T_v} t_i
<
\min_{j\in\mathcal V_v} t_j,
\qquad v=1,\ldots,7.
\end{equation}
}
\textcolor{impGreen}{Thus, each validation block evaluates predictions on sales occurring after the observations used to fit that fold's model. The training window expands across folds. All training-dependent quantities used by the proposed objectives, including $\bar y$ and the centered log-price deviations, are recomputed from $\mathcal T_v$ only.}

\textcolor{impGreen}{\todo{Insert the exact per-fold training/validation counts from the current fold output; they are not available from the aggregate run summary.}}

\begin{table}[!ht]
\centering
% \color{impGreen}
\caption{Rolling-origin validation structure for the final CCAO development sample.}
\label{tab:fold_structure}
\begin{tabular}{cllrr}
\toprule
Fold & Training period & Validation period & $n_{\mathrm{train}}$ & $n_{\mathrm{val}}$ \\
\midrule
1 & 2016-01-01--2017-04-01 (oldest 90\%) & Same window (newest 10\%) & 46,888 & 5,209 \\
2 & 2016-01-01--2018-07-01 (oldest 90\%) & Same window (newest 10\%) & 100,776 & 11,197 \\
3 & 2016-01-01--2019-10-01 (oldest 90\%) & Same window (newest 10\%) & 151,187 & 16,798 \\
4 & 2016-01-01--2021-01-01 (oldest 90\%) & Same window (newest 10\%) & 200,908 & 22,323 \\
5 & 2016-01-01--2022-04-01 (oldest 90\%) & Same window (newest 10\%) & 252,487 & 28,054 \\
6 & 2016-01-01--2023-07-01 (oldest 90\%) & Same window (newest 10\%) & 298,024 & 33,113 \\
7 & 2016-01-01--2023-11-09 (oldest 90\%) & Same window (newest 10\%) & 310,149 & 34,461 \\
\bottomrule
\end{tabular}
\end{table}

The penalty parameter $\rho$ governs an accuracy--equity tradeoff whose appropriate value cannot be determined from the training objective alone. A very small value may leave the baseline trend largely unchanged, whereas a very large value may distort predictions or overfit the equity pattern in the training period. We therefore treat $\rho$ as a model-selection parameter and choose it using chronological validation.

\textcolor{impGreen}{The candidate set contains the standard LightGBM baseline, a custom-objective implementation control with $\rho=0$, and the direct and sample-additive covariance-penalized LightGBM families. For each positive-penalty family, the current experimental specification evaluates 50 geometrically spaced values between $0.1$ and $100$. Because the direct and surrogate objectives have different numerical scales, their $\rho$ values are tuned separately and are not interpreted as comparable magnitudes. The base LightGBM hyperparameters are frozen before this $\rho$ sweep and held fixed across the LightGBM variants.}

% \textcolor{blue}{\sout{Let $\mathcal C$ be a finite set of candidate configurations containing the unpenalized baseline and a prespecified grid of positive penalty values. The candidate set may include both the direct and surrogate formulations. Let $\mathcal V$ index the temporal validation folds, and let $\alpha_v\ge 0$ denote prespecified fold weights satisfying $\sum_{v\in\mathcal V}\alpha_v=1$. For candidate $\eta\in\mathcal C$, let $\operatorname{Perf}_v(\eta)$ be a predictive-performance score on fold $v$, oriented so that larger values are preferred. Let $Q_{m,v}(\eta)$ be the value of ratio-study statistic $m$ on that fold, and let $[\underline\kappa_m,\overline\kappa_m]$ be its prespecified validation range. If $\mathcal M$ denotes the set of active safeguards, the default selection rule is}}

% \begin{center}
% \textcolor{blue}{\sout{\mbox{$\displaystyle
% \eta^\star
% \in
% \argmaxop_{\eta\in\mathcal C}
% \sum_{v\in\mathcal V}\alpha_v\operatorname{Perf}_v(\eta)
% \quad\text{s.t.}\quad
% \underline\kappa_m
% \le
% \sum_{v\in\mathcal V}\alpha_v Q_{m,v}(\eta)
% \le
% \overline\kappa_m,
% \;m\in\mathcal M.
% $}}}
% \end{center}

\textcolor{impGreen}{For the principal CCAO specification, fold-level quantities are averaged with equal weight. For candidate $\eta$, define}
{\color{impGreen}
\begin{equation}
\label{eq:ccao_validation_averages}
\overline{R^2_P}(\eta)
=
\frac{1}{7}\sum_{v=1}^{7}R^2_{P,v}(\eta),
\qquad
\overline{\operatorname{PRD}}(\eta)
=
\frac{1}{7}\sum_{v=1}^{7}\operatorname{PRD}_v(\eta).
\end{equation}
}
\textcolor{impGreen}{The displayed CCAO model is selected by the assessor-oriented rule}
{\color{impGreen}
\begin{equation}
\label{eq:ccao_selection}
\eta^\star
\in
\argmaxop_{\eta\in\mathcal C}
\overline{R^2_P}(\eta)
\qquad
\text{subject to}
\qquad
0.98
\le
\overline{\operatorname{PRD}}(\eta)
\le
1.03.
\end{equation}
}
\textcolor{impGreen}{Thus, the active validation safeguard is PRD, while predictive performance determines the ranking among feasible candidates. PRB, VEI, assessment level, COD, COV, and ratio curves are retained as separate diagnostics rather than folded into the selection objective. This distinction prevents the model-selection rule from collapsing unlike assessment criteria into an arbitrary weighted score.}

\textcolor{blue}{\sout{The averaged constraints in~\eqref{eq:selection} allow a temporary deviation in one validation period to be offset by stronger performance in another. When stability across reassessment periods is a primary concern, a stricter rule can instead require foldwise safeguards. The choice between average-fold and every-fold safeguards should be made before the final test period is examined. Likewise, the candidate grid, active statistics, reference ranges, fold weights, and any tie-breaking rule should be prespecified.}}

\textcolor{impGreen}{Because the rule uses average fold performance, a deviation in one validation period can be offset by another period. We therefore report the fold-level diagnostics of the selected candidate in addition to their average; alternative worst-fold safeguards are treated as sensitivity analyses rather than as part of the headline selection rule. If no candidate satisfies the PRD safeguard, the procedure returns no feasible model rather than relaxing the criterion after observing the evaluation period.}

\textcolor{blue}{\sout{The selection procedure deliberately distinguishes between statistics used as constraints and statistics reported for evaluation. An office may, for example, use PRD as an active validation safeguard while reporting PRB, VEI, COD, assessment level, and subgroup ratio curves as additional diagnostics. The empirical design in the next section specifies the active safeguards used in the CCAO application.}}

\textcolor{blue}{\sout{If no candidate satisfies the prespecified constraints, formulation~\eqref{eq:selection} returns no feasible model. This outcome should prompt examination of the data, base learner, candidate set, and source of the conflicting diagnostics rather than an unreported relaxation of the safeguards. Pareto-based rules and convex stacking can be useful for exploratory analysis of such cases, but they introduce additional choices and are therefore treated as secondary extensions in Appendix~\ref{app:selection_extensions}. The final test period remains untouched until the complete validation-based selection procedure has been applied.}}

\textcolor{impGreen}{After $\eta^\star$ is selected, that configuration is refit on the full 344,610-sale development pool and evaluated on the 38,290-sale held-out test without further tuning. For the 2025 evaluation, the selected configuration is then refit on all 382,900 eligible 2016--2024 sales; the 26,641 2025 sales are used only for this separate forward evaluation and do not trigger re-selection.}

\textcolor{impGreen}{All scalar evaluation statistics are computed on the full relevant validation or evaluation sample after predictions are transformed back to the price scale. For the descriptive ratio curves used in the Results, we plot $r_i=\widehat P_i/P_i$ against $\log P_i$ and overlay a LOWESS curve with smoothing fraction $0.4$. A random subsample of 3,000 observations may be used to render the scatter layer for readability, but the fitted models and reported scalar metrics use the full evaluation sample.}

\begin{table}[t]
\centering
\caption{CCAO application and evaluation design.}
\label{tab:ccao_design}
\begin{tabular}{ll}
\toprule
Item & Design \\
\midrule
Population & Verified Cook County residential sales \\
\textcolor{blue}{\sout{Development period}} & \textcolor{blue}{\sout{2016--2022}} \\
\textcolor{impGreen}{Development / CV} & \textcolor{impGreen}{Oldest 90\% of eligible 2016--2024 sales: 344,610, through 2023-11-09} \\
\textcolor{blue}{\sout{Untouched test period}} & \textcolor{blue}{\sout{2023}} \\
\textcolor{impGreen}{Primary held-out test} & \textcolor{impGreen}{Newest 10\%: 38,290 sales, 2023-11-09--2024-12-31} \\
\textcolor{impGreen}{Secondary forward evaluation} & \textcolor{impGreen}{26,641 sales, 2025-01-01--2025-12-29} \\
\textcolor{blue}{\sout{Pre-2024 observations}} & \textcolor{blue}{\sout{349,661}} \\
\textcolor{impGreen}{Production training} & \textcolor{impGreen}{382,900 eligible 2016--2024 sales} \\
Predictors & 95 total; 23 categorical \\
Validation & Seven expanding-window rolling-origin folds; 15-month steps; newest 10\% validates \\
Base nonlinear learner & LightGBM on log sale price \\
\textcolor{blue}{\sout{Primary correction}} & \textcolor{blue}{\sout{Covariance-guided sample-additive surrogate}} \\
\textcolor{impGreen}{Corrections evaluated} & \textcolor{impGreen}{Direct squared covariance and sample-additive covariance upper bound} \\
\textcolor{blue}{\sout{Selection}} & \textcolor{blue}{\sout{Best validation accuracy among candidates meeting active safeguards}} \\
\textcolor{impGreen}{Principal selection rule} & \textcolor{impGreen}{Highest mean validation $R^2_P$ among candidates with mean PRD in $[0.98,1.03]$} \\
\bottomrule
\end{tabular}
\end{table}

\subsection{\textcolor{blue}{\sout{Required Comparators}}\textcolor{impGreen}{Models and Comparators}}
\label{subsec:comparators}

\textcolor{blue}{\sout{Four comparisons are needed to interpret the method.}}
\textcolor{impGreen}{The empirical comparison separates the baseline model class, the effect of entering the custom-objective implementation, the effect of positive covariance regularization, and simpler post-hoc correction.}

\paragraph{Linear and nonlinear baselines.}
\textcolor{blue}{\sout{Linear regression shows the performance of a simpler model; unpenalized LightGBM is the relevant modern AVM baseline.}}
\textcolor{impGreen}{Linear regression provides the simpler historically motivated benchmark; standard LightGBM, trained with the ordinary squared log-price objective, is the relevant nonlinear AVM baseline. Both use the same outcome, development observations, and underlying predictor information under the common temporal design.}

\paragraph{Direct and surrogate penalties.}
The direct penalty tests the stated covariance target. The surrogate tests the implementation recommended for standard boosting software.
\textcolor{blue}{\sout{A custom-objective fit at $\rho=0$ is also needed because LightGBM initialization and training behavior can differ between an ordinary regression objective and a user-supplied objective. This control separates the effect of the penalty from the effect of entering the custom-objective code path.}}
\textcolor{impGreen}{The direct LightGBM comparator uses the diagonal-Hessian approximation described in Appendix~\ref{app:implementation}: its covariance gradient is exact, while the cross-observation second-order curvature is omitted so that the objective can be used through the standard LightGBM interface. The surrogate uses the exact observation-additive objective from Section~\ref{subsec:surrogate}. The final comparison also includes a custom-objective fit at $\rho=0$. This implementation control separates changes caused by the custom-objective code path from changes caused by a positive penalty.}

\paragraph{Post-hoc rescaling.}
\textcolor{blue}{\sout{A one-parameter rescaling fitted on training or validation data is a low-cost competitor. It must be reported because a global correction may reproduce part of the penalty's effect without changing the training objective. The paper therefore does not treat improvement over unpenalized LightGBM as sufficient evidence that custom regularization is necessary.}}
\textcolor{impGreen}{A validation-fitted one-parameter rescaling of the baseline predictions is included as a low-cost post-hoc competitor where reported in the Results, particularly in the multi-county benchmark. This comparison is necessary because a global spread correction may reproduce part of the penalty's effect without modifying model training; improvement relative to standard LightGBM alone is therefore not evidence that in-processing regularization is necessary.}

A constant multiplicative level adjustment and a slope-changing rescaling serve different roles. Multiplying every price prediction by a constant adds a constant to every log residual and therefore leaves $\operatorname{Cov}(e,y)$ unchanged. By contrast, a one-degree calibration that changes the spread of predictions can alter the ratio--value slope and is the relevant competitor in the multi-county benchmark.

\textcolor{impGreen}{\todo{Insert the exact centered one-parameter recalibration map and its validation fitting criterion here or in Appendix~\ref{app:attom}; the final manuscript should not leave this comparator defined only verbally.}}

\paragraph{\textcolor{blue}{\sout{Richer information.}}}
\textcolor{blue}{\sout{Spatial, temporal, neighborhood, and property-characteristic improvements can reduce omitted structure and may improve both accuracy and equity. They are complements to the penalty, not substitutes assumed away by the paper. The preferred sequence is to build the strongest defensible AVM first and then test whether an explicit vertical-equity safeguard still adds value.}}

\textcolor{impGreen}{Better spatial, temporal, neighborhood, and property information can improve both prediction and equity, but this is a modeling extension rather than a comparator in the headline experiment. We therefore discuss richer information as a complementary direction in Section~\ref{sec:discussion} rather than treating it as a completed empirical control here.}


\subsection{Exploratory Six-County ATTOM Benchmark}
\label{subsec:attom_design}

We also run a common ATTOM-derived pipeline in Cook County, IL; Allegheny County, PA; Maricopa County, AZ; King County, WA; Miami-Dade County, FL; and Middlesex County, MA. The benchmark contains 1,231,194 qualified single-family sales and 246,241 held-out sales. Each county uses chronological validation and test blocks and the same modeling code path.

\textcolor{blue}{\sout{County-specific penalty values are calibrated to the baseline prediction scale.}}
\textcolor{impGreen}{For each county, eligible recorder transfers are joined to the most recent assessor-history record available before the sale, together with lagged contextual information. The benchmark then orders eligible sales chronologically and reserves the final 20\% as its test block, with model and penalty selection performed only on earlier observations. The benchmark compares standard LightGBM, the custom-objective $\rho=0$ control, positive covariance-penalty values, and validation-fitted centered recalibration of the baseline log predictions. The exact county periods, sample counts, common filters, feature set, and recalibration specification are reported in Appendix~\ref{app:attom}.}

\textcolor{impGreen}{\todo{Add a compact county-by-county design table here if these dates/counts are needed to read the main Results without consulting the appendix. At minimum, archive them in Appendix~\ref{app:attom}.}}

This is an exploratory transfer benchmark, not an audit of the six assessor offices. The data are commercially harmonized and do not reproduce each jurisdiction's sales-verification rules, local features, legal setting, or production review. The benchmark was developed iteratively, and its present county configurations were not all frozen under one confirmatory protocol before test evaluation. We therefore use it to examine direction, heterogeneity, and failure modes. It does not provide confirmatory estimates of performance in the six jurisdictions.

\textcolor{impGreen}{Accordingly, the six-county results are kept separate from the primary CCAO application: they test whether the qualitative behavior of the correction transfers across heterogeneous datasets, not whether the six offices would obtain the reported results in their own production systems.}

\section{\blueheading{Application and Evaluation Design}}
\label{sec:empirical_design}

\subsection{\blueheading{Application to the Translated CCAO Residential Pipeline}}
\label{subsec:ccao_design}
{\color{red} For the empirical analysis, we implement in Python the components of CCAO's residential valuation workflow required for model development and evaluation. The implementation closely follows CCAO's data-preparation, model-fitting, and ratio-study procedures, preserving the operational structure of the existing workflow while enabling controlled evaluation of alternative training objectives.

The proposed methodology modifies only the model-training objective. We hold fixed the sale filters, predictor set, temporal data splits, transformation of predictions to the price scale, ratio-study procedures, and surrounding review process. This restricted intervention serves two purposes. First, it isolates the effect of incorporating vertical equity into model estimation. Second, it demonstrates how the methodology can be integrated into an existing mass-appraisal workflow without redesigning the broader valuation system.}
\advisorchange{The CCAO application uses verified Cook County residential sales} from the CCAO 2025 assessment-cycle data. Following the office pipeline, we exclude multicard PINs and sales marked as outliers. The sales universe is 2016--2024 and contains 382,900 eligible observations; the current file contains 2016 sales, so observed sales begin on 2016-01-01. The oldest 90\% (344,610 sales, through 2023-11-09) form the development/CV pool; the newest 10\% (38,290 sales, 2023-11-09 through 2024-12-31) form the held-out test. The 2025 evaluation contains 26,641 sales and is run after production refitting on all 382,900 eligible 2016--2024 sales.

The model uses 95 predictors, including property characteristics, location and proximity variables, Census measures, time variables, parcel-shape measures, neighborhood identifiers, and other assessor variables. Twenty-three are treated as categorical. All predictors are constructed upstream; the penalized models use the same feature set as the LightGBM baseline.

Within the 344,610-sale development/CV pool, we use seven expanding-window rolling-origin folds stepping every 15 months from the 2016-01-01 anchor; in each cumulative window, the newest 10\% of observations form validation. For the principal CCAO comparison, candidates are screened using the average validation PRD reference band and then ranked by predictive performance. PRB, VEI, assessment level, and dispersion are reported as additional diagnostics rather than used to select the displayed model. The 38,290-sale held-out test is evaluated only after selection.

\begin{table}[t]
\centering
\caption{\advisorchange{CCAO application and evaluation design.}}
\label{tab:ccao_design}
\begin{tabular}{ll}
\toprule
Item & Design \\
\midrule
Population & Verified Cook County residential sales \\
Development / CV & Oldest 90\% of eligible 2016--2024 sales: 344,610 through 2023-11-09 \\
Held-out test & Newest 10\%: 38,290 sales, 2023-11-09--2024-12-31 \\
Production training & 382,900 eligible 2016--2024 sales \\
2025 forward evaluation & 26,641 sales \\
Predictors & 95 total; 23 categorical \\
Validation & Seven expanding-window rolling-origin folds; 15-month steps; newest 10\% validates \\
Base nonlinear learner & LightGBM on log sale price \\
Primary correction & Covariance-guided sample-additive surrogate \\
Selection & Best validation accuracy among candidates meeting active safeguards \\
\bottomrule
\end{tabular}
\end{table}

\subsection{Required Comparators}
\label{subsec:comparators}

Four comparisons are needed to interpret the method.

\paragraph{Linear and nonlinear baselines.}
Linear regression shows the performance of a simpler model; unpenalized LightGBM is the relevant modern AVM baseline.

\paragraph{Direct and surrogate penalties.}
The direct penalty tests the stated covariance target. The surrogate tests the implementation recommended for standard boosting software. A custom-objective fit at $\rho=0$ is also needed because LightGBM initialization and training behavior can differ between an ordinary regression objective and a user-supplied objective. This control separates the effect of the penalty from the effect of entering the custom-objective code path.

\paragraph{Post-hoc rescaling.}
A one-parameter rescaling fitted on training or validation data is a low-cost competitor. It must be reported because a global correction may reproduce part of the penalty's effect without changing the training objective. The paper therefore does not treat improvement over unpenalized LightGBM as sufficient evidence that custom regularization is necessary.

A constant multiplicative level adjustment and a slope-changing rescaling serve different roles. Multiplying every price prediction by a constant adds a constant to every log residual and therefore leaves $\Cov(e,y)$ unchanged. By contrast, a one-degree calibration that changes the spread of predictions can alter the ratio--value slope and is the relevant competitor in the multi-county benchmark.

\paragraph{Richer information.}
Spatial, temporal, neighborhood, and property-characteristic improvements can reduce omitted structure and may improve both accuracy and equity. They are complements to the penalty, not substitutes assumed away by the paper. The preferred sequence is to build the strongest defensible AVM first and then test whether an explicit vertical-equity safeguard still adds value.

\subsection{\blueheading{Exploratory Six-County ATTOM Benchmark}}
\label{subsec:attom_design}

We also run a common ATTOM-derived pipeline in Cook County, IL; Allegheny County, PA; Maricopa County, AZ; King County, WA; Miami-Dade County, FL; and Middlesex County, MA. The benchmark contains 1,231,194 qualified single-family sales and 246,241 held-out sales. Each county uses chronological validation and test blocks, the same modeling code path, and county-specific penalty values calibrated to the baseline prediction scale.

\advisorchange{This is an exploratory transfer benchmark, not an audit of the six assessor offices.} The data are commercially harmonized and do not reproduce each jurisdiction's sales-verification rules, local features, legal setting, or production review. The benchmark was developed iteratively, and its present county configurations were not all frozen under one confirmatory protocol before test evaluation. We therefore use it to examine direction, heterogeneity, and failure modes. It does not provide confirmatory estimates of performance in the six jurisdictions.


{
\color{red}
\subsection{Validation-Based Model Selection}
\label{subsec:selection}

\todo{Either correctly connect this or move to appendix, and just reference it}

The penalty parameter \(\rho\) governs an accuracy--equity tradeoff whose appropriate value cannot be determined from the training objective alone. A very small value may leave the baseline trend largely unchanged, whereas a very large value may distort predictions or overfit the equity pattern in the training period. We therefore treat \(\rho\) as a model-selection parameter and choose it using chronological validation.

Let \(\mathcal C\) be a finite set of candidate configurations containing the unpenalized baseline and a prespecified grid of positive penalty values. The candidate set may include both the direct and surrogate formulations. Let \(\mathcal V\) index the temporal validation folds, and let \(\alpha_v\ge 0\) denote prespecified fold weights satisfying
\[
\sum_{v\in\mathcal V}\alpha_v=1.
\]
For candidate \(\eta\in\mathcal C\), let \(\operatorname{Perf}_v(\eta)\) be a predictive-performance score on fold \(v\), oriented so that larger values are preferred. Let \(Q_{m,v}(\eta)\) be the value of ratio-study statistic \(m\) on that fold, and let \([\underline\kappa_m,\overline\kappa_m]\) be its prespecified validation range. If \(\mathcal M\) denotes the set of active safeguards, the default selection rule is
\begin{align}
\label{eq:selection}
\eta^\star
\in
\argmaxop_{\eta\in\mathcal C}
\quad
& \sum_{v\in\mathcal V}
\alpha_v\operatorname{Perf}_v(\eta)
\\
\text{s.t.}\quad
&
\underline\kappa_m
\le
\sum_{v\in\mathcal V}
\alpha_v Q_{m,v}(\eta)
\le
\overline\kappa_m,
\qquad
m\in\mathcal M.
\nonumber
\end{align}
This formulation mirrors the operational decision problem: among models satisfying the active ratio-study safeguards, select the one with the strongest out-of-time predictive performance. It avoids collapsing prediction, assessment level, dispersion, and vertical equity into a single weighted score whose interpretation would depend on arbitrary scaling choices.

The averaged constraints in~\eqref{eq:selection} allow a temporary deviation in one validation period to be offset by stronger performance in another. When stability across reassessment periods is a primary concern, a stricter rule can instead require
\[
\underline\kappa_m
\le
Q_{m,v}(\eta)
\le
\overline\kappa_m,
\qquad
m\in\mathcal M,\quad v\in\mathcal V.
\]
The choice between average-fold and every-fold safeguards should be made before the final test period is examined. Likewise, the candidate grid, active statistics, reference ranges, fold weights, and any tie-breaking rule should be prespecified.

The selection procedure deliberately distinguishes between statistics used as constraints and statistics reported for evaluation. An office may, for example, use PRD as an active validation safeguard while reporting PRB, VEI, COD, assessment level, and subgroup ratio curves as additional diagnostics. The empirical design in the next section specifies the active safeguards used in the CCAO application.

If no candidate satisfies the prespecified constraints, formulation~\eqref{eq:selection} returns no feasible model. This outcome should prompt examination of the data, base learner, candidate set, and source of the conflicting diagnostics rather than an unreported relaxation of the safeguards. Pareto-based rules and convex stacking can be useful for exploratory analysis of such cases, but they introduce additional choices and are therefore treated as secondary extensions in Appendix~\ref{app:selection_extensions}. The final test period remains untouched until the complete validation-based selection procedure has been applied.




}







\section{Results}
\label{sec:results}

\subsection{Accuracy Alone Does Not Remove the CCAO Ratio Trend}
\label{subsec:ccao_baseline_results}

Table~\ref{tab:ccao_main_results} reports the principal 2023 test-set comparison. LightGBM improves $R^2$ from $0.762$ to $0.885$ and lowers MAE from about \$86,244 to \$71,912 relative to linear regression. COD also falls. The vertical-equity measures move in the opposite direction: PRD rises to $1.085$, PRB falls to $-0.118$, and VEI falls to $-37.64\%$.

The validation-selected surrogate model uses $\rho=5.96$. Relative to the LightGBM baseline, it reduces the PRD deviation from one from $0.085$ to $0.039$, the absolute PRB from $0.118$ to $0.014$, and the absolute VEI from $37.64\%$ to $9.27\%$. $R^2$ decreases by $0.004$, and MAE increases by about $3.5\%$. COD and COV improve modestly. The test-set PRD point estimate remains above the $1.03$ reference bound. The distinction is substantive: the selected model improves the ratio pattern but does not satisfy every test-set reference band. The current CCAO results table does not include a custom-objective zero-penalty control. The comparison therefore measures the total change from the ordinary LightGBM fit to the implemented penalized fit; it does not isolate the penalty from every difference introduced by the custom-objective training path.

\begin{table}[t]
\centering
\caption{Principal CCAO 2023 test-set results. The penalized model was selected using rolling-origin validation, not the test set.}
\label{tab:ccao_main_results}
\resizebox{\textwidth}{!}{%
\begin{tabular}{lrrrrrrrrrr}
\toprule
Model & $\rho$ & $R^2$ & MAE & Median $r$ & Mean $r$ & COD & COV & PRD & PRB & VEI \\
\midrule
Linear regression & -- & 0.762 & \$86,244 & 1.001 & 1.050 & 25.01\% & 36.86\% & 1.039 & -0.011 & -13.53\% \\
LightGBM baseline & -- & 0.885 & \$71,912 & 0.944 & 1.021 & 23.07\% & 35.48\% & 1.085 & -0.118 & -37.64\% \\
Validation-selected surrogate & 5.96 & 0.881 & \$74,416 & 0.937 & 0.978 & 22.80\% & 33.64\% & 1.039 & -0.014 & -9.27\% \\
\bottomrule
\end{tabular}}
\end{table}

\subsection{The Correction Acts in the Intended Direction}
\label{subsec:mechanism_results}

Figure~\ref{fig:penalty_mechanism} compares the baseline ratio curve with representative direct and surrogate corrections. The baseline curve declines sharply with value. Both penalized fits flatten the global trend, with the surrogate producing the largest reduction in this CCAO run. These plots are mechanism checks, not substitutes for PRD, PRB, VEI, assessment-level statistics, or subgroup analysis.

\begin{figure}[H]
\centering
\begin{subfigure}{0.48\textwidth}
\centering
\safeincludegraphics[width=\linewidth]{slides/img/Baseline LightGBM_0_motivation.pdf}
\caption{Baseline LightGBM}
\end{subfigure}\hfill
\begin{subfigure}{0.48\textwidth}
\centering
\safeincludegraphics[width=\linewidth]{slides/img/Covariance-Penalized LightGBM_2_motivation.pdf}
\caption{Direct squared-covariance penalty}
\end{subfigure}

\vspace{3mm}
\begin{subfigure}{0.48\textwidth}
\centering
\safeincludegraphics[width=\linewidth]{slides/img/Covariance-Penalized LightGBM_3_motivation.pdf}
\caption{Covariance-guided surrogate}
\end{subfigure}
\caption{Assessment-ratio patterns in the translated CCAO test set. The penalty is intended to flatten the global value-related trend; all assessor-facing measures remain necessary for evaluation.}
\label{fig:penalty_mechanism}
\end{figure}


{
\color{red}
\begin{table}[!htbp]
\centering
% \caption{Held-out test-set performance for representative penalized LightGBM configurations. The penalty parameter $\rho$ is penalty-specific and should not be compared across different penalty definitions.}
\caption{PLACEHOLDER OF THE 3 REPRESENTATIVE PENALIZED RESULTS}
\label{tab:penalized_models}
\resizebox{\textwidth}{!}{%
\begin{tabular}{lc|cc|cc|cc|ccc}
\toprule
Model Family & $\rho$ & $R^2$ & MAE & Med.\ $r$ & Mean $r$ & COD & COV & VEI & PRD & PRB \\
\midrule
Linear Regression  & -- & 0.762 & \$86{,}256.879 & 1.001 & 1.051 & 25.013\% & 36.825\% & -13.474\% & 1.039 & -0.011 \\
LGBMRegressor      & -- & 0.883 & \$72{,}063.853 & 0.944 & 1.021 & 23.106\% & 35.398\% & -36.816\% & 1.085 & -0.118 \\
\midrule
LGBSurrogatePenalty & 0.100 & 0.881 & \$73{,}271.882 & 0.936 & 1.011 & 23.125\% & 35.236\% & -36.367\% & 1.084 & -0.115 \\
LGBSurrogatePenalty & 0.189 & 0.881 & \$73{,}208.883 & 0.936 & 1.010 & 23.042\% & 35.048\% & -33.608\% & 1.080 & -0.107 \\
LGBSurrogatePenalty & 0.404 & 0.882 & \$73{,}110.847 & 0.937 & 1.006 & 22.924\% & 34.782\% & -29.878\% & 1.075 & -0.095 \\
LGBSurrogatePenalty & 0.761 & 0.883 & \$72{,}974.386 & 0.937 & 1.001 & 22.749\% & 34.189\% & -23.789\% & 1.067 & -0.076 \\
LGBSurrogatePenalty & 1.628 & 0.882 & \$73{,}233.014 & 0.938 & 0.994 & 22.645\% & 33.664\% & -16.431\% & 1.057 & -0.052 \\
LGBSurrogatePenalty & 3.070 & 0.881 & \$73{,}663.987 & 0.937 & 0.986 & 22.667\% & 33.620\% & -12.630\% & 1.048 & -0.032 \\
LGBSurrogatePenalty & 6.572 & 0.882 & \$74{,}428.114 & 0.937 & 0.976 & 22.806\% & 33.386\% & -8.545\% & 1.038 & -0.011 \\
LGBSurrogatePenalty & 12.390 & 0.879 & \$75{,}846.355 & 0.935 & 0.968 & 23.249\% & 33.684\% & -7.446\% & 1.032 & 0.003 \\
% LGBSurrogatePenalty & 26.519 & 0.879 & \$77{,}113.783 & 0.933 & 0.960 & 23.719\% & 33.959\% & -4.484\% & 1.024 & 0.020 \\
% LGBSurrogatePenalty & 50.000 & 0.878 & \$78{,}155.478 & 0.932 & 0.956 & 24.092\% & 34.336\% & -4.883\% & 1.021 & 0.027 \\
\midrule
LGBCovPenalty & 0.100 & 0.879 & \$73{,}459.227 & 0.937 & 1.014 & 23.210\% & 35.551\% & -38.368\% & 1.087 & -0.122 \\
LGBCovPenalty & 0.189 & 0.880 & \$73{,}311.803 & 0.939 & 1.015 & 23.221\% & 35.443\% & -37.671\% & 1.086 & -0.119 \\
LGBCovPenalty & 0.404 & 0.882 & \$72{,}947.453 & 0.941 & 1.017 & 23.224\% & 35.432\% & -36.501\% & 1.083 & -0.115 \\
LGBCovPenalty & 0.761 & 0.883 & \$72{,}607.948 & 0.944 & 1.019 & 23.157\% & 35.246\% & -34.017\% & 1.078 & -0.106 \\
LGBCovPenalty & 1.628 & 0.886 & \$72{,}198.279 & 0.949 & 1.023 & 23.176\% & 35.086\% & -30.427\% & 1.071 & -0.096 \\
LGBCovPenalty & 3.070 & 0.884 & \$72{,}293.558 & 0.955 & 1.027 & 23.216\% & 35.053\% & -26.850\% & 1.065 & -0.085 \\
LGBCovPenalty & 6.572 & 0.881 & \$72{,}871.805 & 0.961 & 1.030 & 23.282\% & 35.040\% & -20.525\% & 1.055 & -0.068 \\
LGBCovPenalty & 12.390 & 0.878 & \$73{,}449.207 & 0.963 & 1.030 & 23.261\% & 34.757\% & -16.050\% & 1.048 & -0.057 \\
LGBCovPenalty & 26.519 & 0.870 & \$74{,}741.535 & 0.963 & 1.028 & 23.310\% & 34.542\% & -13.572\% & 1.041 & -0.046 \\
LGBCovPenalty & 50.000 & 0.866 & \$75{,}574.792 & 0.960 & 1.023 & 23.357\% & 34.522\% & -11.105\% & 1.037 & -0.039 \\
\bottomrule
\end{tabular}}
\end{table}

}

\subsection{The Penalty Traces a Usable Validation Path}
\label{subsec:tradeoff_results}

As $\rho$ increases, the candidate models move from the accuracy-only LightGBM solution toward flatter ratio curves. The change is not a fixed or necessarily monotone tradeoff across all metrics. A moderate range gives large improvements in PRB and VEI with small changes in $R^2$; larger penalties eventually increase predictive error and can move the model toward progressivity. This pattern supports a grid search with temporal validation rather than a single theoretical value of $\rho$.

\begin{figure}[H]
\centering
\begin{subfigure}{0.48\textwidth}
\centering
\safeincludegraphics[width=\textwidth]{slides/img/tradeoffs/PRD_vs_R2_band.pdf}
\caption{PRD versus $R^2$}
\end{subfigure}\hfill
\begin{subfigure}{0.48\textwidth}
\centering
\safeincludegraphics[width=\textwidth]{slides/img/tradeoffs/VEI_vs_R2_band.pdf}
\caption{VEI versus $R^2$}
\end{subfigure}
\caption{CCAO accuracy--equity paths across penalty values. Reference bands guide candidate screening; the final test period is not used to choose $\rho$.}
\label{fig:tradeoffs}
\end{figure}

The exact regularization path derived in Appendix~\ref{app:theory} explains why moderate penalties can remove a first-order component at second-order in-sample squared-error cost within a fixed prediction space. For fully retrained boosted trees, however, split structure and early stopping can change. The theory is used to interpret and scale the grid, \advisorchange{not to replace out-of-time validation}.

\subsection{Exploratory Six-County Evidence Shows Heterogeneity and Limits}
\label{subsec:attom_results}

Table~\ref{tab:attom_baselines} shows the unpenalized ATTOM benchmark results. Every county has at least one vertical-equity point estimate outside the reference bands used in the benchmark, but the patterns are not identical. Cook, Allegheny, Maricopa, and Middlesex have negative PRB values outside the benchmark band. King and Miami-Dade have PRB values within the band even though their PRD values are high. Miami-Dade is especially useful as a boundary case: its positive PRB and curved ratio pattern show why a global covariance target and a single ratio statistic cannot be treated as complete descriptions of vertical equity.

\begin{table}[t]
\centering
\caption{Unpenalized LightGBM results in the six-county ATTOM benchmark. These are held-out model-benchmark point estimates, not official assessor results.}
\label{tab:attom_baselines}
\resizebox{\textwidth}{!}{%
\begin{tabular}{lrrrrrrr}
\toprule
County & $R^2$ (log) & RMSE (log) & MAPE & COD & PRD & PRB & Median $r$ \\
\midrule
Cook, IL & 0.833 & 0.283 & 0.207 & 21.199 & 1.083 & -0.098 & 0.959 \\
Allegheny, PA & 0.779 & 0.325 & 0.253 & 26.129 & 1.103 & -0.101 & 0.954 \\
Maricopa, AZ & 0.844 & 0.222 & 0.147 & 14.797 & 1.072 & -0.055 & 0.989 \\
King, WA & 0.798 & 0.261 & 0.207 & 16.336 & 1.073 & -0.028 & 1.106 \\
Miami-Dade, FL & 0.808 & 0.316 & 0.263 & 24.210 & 1.092 & 0.025 & 1.060 \\
Middlesex, MA & 0.780 & 0.264 & 0.191 & 18.565 & 1.085 & -0.096 & 1.023 \\
\bottomrule
\end{tabular}}
\end{table}

Among the four counties whose baseline PRB point estimate is outside the benchmark band, at least one tested penalty places PRB inside the band in all four, with no more than a 1.84\% increase in log-scale RMSE. In contrast, only Maricopa has a tested penalty that places both PRD and PRB point estimates inside their bands. These descriptive results suggest that the first-order correction can transfer across datasets, but they do not establish a universal or confirmatory compliance result.

\subsection{The Current Evidence Does Not Establish Dominance over Recalibration}
\label{subsec:scaling_results}

\advisorchange{The multi-county benchmark includes validation-fitted centered recalibration maps} of the baseline log predictions. These comparators can change the spread of predictions after model fitting and can therefore alter a global ratio--value trend. A constant multiplicative level adjustment is different: it shifts every log residual by the same amount and leaves the targeted covariance unchanged.

The existing comparisons do not support a general claim that the custom objective dominates the simpler recalibration. The in-processing penalty may still be useful when the target must influence tree construction or when the method will be extended to geographic, class-specific, or other structured constraints. For the current paper, the defensible rule is simpler: report both approaches, fit both without test data, and prefer recalibration when its out-of-time accuracy and equity results are materially the same.

\subsection{Assessor Implementation}
\label{sec:implementation_workflow}

The method can be inserted into an existing AVM with the following sequence.

\begin{enumerate}[leftmargin=7mm]
    \item Build the strongest defensible baseline model using verified sales, appropriate time adjustment, relevant property features, and local market information.
    \item Hold out a future period and create rolling or blocked validation folds that respect sale dates.
    \item Report assessment level, COD, PRD, PRB, supplementary vertical-equity measures, and ratio curves for the baseline model.
    \item Fit the covariance-guided objective over a logarithmic grid of $\rho$, including $\rho=0$ as an implementation control.
    \item Fit a simple post-hoc rescaling benchmark using development data only.
    \item Select the most accurate candidate that meets prespecified validation safeguards, then evaluate it once on the untouched test period.
    \item Repeat the diagnostics by geography, property class, value range, and time. Reject a global improvement that hides material local deterioration.
\end{enumerate}

\begin{table}[t]
\centering
\caption{Different pipeline components address different problems.}
\label{tab:roles}
\begin{tabular}{p{4.2cm}p{10.0cm}}
\toprule
Component & Primary role \\
\midrule
Better property, spatial, and temporal data & Reduce omitted structure and improve the conditional valuation model; may improve accuracy and equity together \\
Constant level calibration & Move the overall ratio level toward its required target; does not change log-residual covariance \\
One-parameter slope/spread rescaling & Low-cost post-hoc correction of a global value trend; required comparator for the penalty \\
Covariance-guided regularization & Discourage a global first-order log-ratio/log-price trend during model training \\
Stratified and local diagnostics & Detect offsetting geographic, class-specific, or nonlinear patterns hidden by aggregate measures \\
Temporal validation and uncertainty analysis & Test whether the selected correction generalizes across market periods \\
\bottomrule
\end{tabular}
\end{table}

The office should save the penalty definition, grid, validation predictions, active constraints, selected configuration, and all ratio-study outputs. This documentation makes the change auditable and allows the correction to be re-evaluated in later assessment cycles.

\section{Discussion and Conclusions}
\label{sec:discussion}

\advisorchange{The CCAO application answers a limited question.} A training-time correction can remove much of a countywide regressive ratio trend while retaining most of the predictive performance of the ordinary LightGBM model. The selected sample-additive model improves all three reported vertical-equity point estimates, although test PRD remains outside its reference band and MAE increases. \advisorchange{The results therefore support use of the penalty as a candidate model}, not as proof that the equity problem has been solved.

The value of the approach is operational. It converts a ratio-study concern into a model-training option, works through a standard boosted-tree interface, and can be tuned with chronological validation. \advisorchange{Those features follow directly from the CCAO problem that motivated the collaboration.}

The evidence also limits the paper's claim. Better data may improve accuracy and equity together. The covariance target is countywide and first-order; it does not set assessment level, guarantee any ratio-study range, or remove local inequity. The current CCAO comparison lacks a custom-objective zero-penalty control, and the exploratory county benchmark does not establish confirmatory external validity. Post-hoc recalibration may reproduce part of the gain. For these reasons, the penalty should enter the same review process as any other model change and should never be treated as an equity certificate.

\subsection{Limitations}

\advisorchange{First, the CCAO application uses sold properties.} Results on sales do not establish performance for the full unsold inventory, particularly when the sale sample is not representative. \advisorchange{Second, the primary application is a Python translation} of the CCAO workflow rather than an official production run. Third, the penalty targets log-scale covariance, while PRD, PRB, and VEI are price-scale statistics with different functional forms. Fourth, the main CCAO comparison does not yet include the custom-objective zero-penalty control needed to isolate implementation effects. Fifth, the main results are point estimates; uncertainty and variation across folds should accompany an operational decision. Sixth, a countywide correction can average away local regressivity and progressivity. Seventh, the ATTOM benchmark is exploratory, uses commercially harmonized data, and does not reproduce the six offices' institutional data or a single confirmatory model-selection protocol. Finally, the current evidence does not show that in-processing regularization is better than post-hoc recalibration.

\subsection{Future Work}

\advisorchange{The next evaluation step is a clean ablation} that compares the baseline model, better spatial and temporal information, the covariance correction, and their combination under the same temporal splits. Comparable-property and neighborhood features must use only information available before the target sale. Preliminary tests of second- and third-degree global recalibration have shown little consistent improvement over a first-degree map, so higher-order corrections should remain secondary unless they add stable out-of-time gains.

The main methodological extension is structured covariance control by geography, property class, and time, with safeguards for small or unstable groups. That extension should be judged against stratified post-hoc recalibration and against a stronger base AVM. Further work should also add block-bootstrap uncertainty, worst-period validation, the custom-objective zero-penalty control, and evaluation on the full assessment inventory. Desk review by assessors is needed to determine whether the changed predictions remain understandable and operationally defensible.




\printbibliography


\appendix

\section{Metric Details and Guidance Status}
\label{app:metrics}

The 2013 IAAO Standard on Ratio Studies is the established standard cited for PRD and COD guidance in this paper \cite{IAAO2013Ratio}. The May 2026 revision of the ratio-studies document is an exposure draft. Measures introduced or revised there, including VEI, are reported as supplementary diagnostics \cite{IAAO2026ExposureRatio}. Reference bands are used for candidate screening; the paper does not infer legal compliance from a point estimate.

The weighted mean ratio satisfies
\[
\bar r_W
=
\frac{\sum_i r_iP_i}{\sum_iP_i}
=
\frac{\sum_i\widehat P_i}{\sum_iP_i},
\]
and PRD is $\bar r/\bar r_W$. PRB follows the IAAO median-normalized regression definition. The full implementation \advisorchange{used in the application should be released} with the code so that all transformations, trimming rules, and proxy definitions are reproducible.

\section{Squared-Error Prediction and Residual--Value Dependence}
\label{app:population_identity}

Let $Y=\log P$, let $X$ be the observed feature vector, and let
\[
f^\star(X)=\E[Y\mid X]
\]
be the squared-error Bayes predictor. Define $e^\star=f^\star(X)-Y$. If $\E[Y^2]<\infty$, then
\begin{proposition}
\label{prop:bayes_covariance}
\[
\Cov(e^\star,Y)
=
-\E[\Var(Y\mid X)]
\le 0.
\]
\end{proposition}
\begin{proof}
Using the law of total variance,
\[
\Cov(e^\star,Y)
=
\Cov(\E[Y\mid X]-Y,Y)
=
\Var(\E[Y\mid X])-\Var(Y)
=
-\E[\Var(Y\mid X)].
\]
\end{proof}

This identity explains a source of negative log-residual dependence when important price variation remains unresolved after conditioning on $X$. It does not imply that every finite-sample model is regressive under every assessor metric. It also does not imply that adding features cannot improve both accuracy and equity.

\section{Covariance and PRD on the Price Scale}
\label{app:cov_prd}

Let $P>0$ and $r=\widehat P/P>0$. The population PRD can be written as
\[
\mathrm{PRD}
=
\frac{\E[r]\E[P]}{\E[rP]}.
\]
Then
\begin{proposition}
\label{prop:cov_prd}
\[
\operatorname{sign}\!\left(\Cov(r,P)\right)
=
\operatorname{sign}(1-\mathrm{PRD}).
\]
\end{proposition}
\begin{proof}
Since $\E[rP]>0$,
\[
1-\mathrm{PRD}
=
\frac{\E[rP]-\E[r]\E[P]}{\E[rP]}
=
\frac{\Cov(r,P)}{\E[rP]}.
\]
\end{proof}

The penalty in the paper acts on $\Cov(\log r,\log P)$, not on $\Cov(r,P)$. Proposition~\ref{prop:cov_prd} therefore gives directional motivation rather than an equality between the training target and PRD.

\section{Exact Fixed-Space Regularization Path}
\label{app:theory}

Let $\Vcal\subseteq\R^n$ be a closed linear prediction space, define
\[
\langle u,v\rangle_n=\frac{1}{n}u^\top v,
\qquad
\|u\|_n^2=\langle u,u\rangle_n,
\qquad
c=y-\bar y\1,
\]
and consider
\begin{equation}
\label{eq:path_problem}
f_\rho\in\argminop_{f\in\Vcal}
\left\{
\|f-y\|_n^2
+\frac{\rho}{2}\langle f-y,c\rangle_n^2
\right\}.
\end{equation}
Let $P_{\Vcal}$ be the orthogonal projection under $\langle\cdot,\cdot\rangle_n$ and set
\[
f_0=P_{\Vcal}y,\qquad
C_0=\langle f_0-y,c\rangle_n,\qquad
g=P_{\Vcal}c,\qquad
A=\|g\|_n^2.
\]

\begin{proposition}[Covariance shrinkage path]
\label{prop:path}
If $A>0$, then
\[
f_\rho=f_0-\frac{\rho}{2}C_\rho g,
\qquad
C_\rho=\frac{C_0}{1+\rho A/2}.
\]
Consequently, the slope of residuals on $y$ is multiplied by
\[
q(\rho)=\frac{1}{1+\rho A/2}.
\]
\end{proposition}

\begin{proof}
For every $h\in\Vcal$, first-order optimality gives
\[
2\langle f_\rho-y,h\rangle_n
+\rho C_\rho\langle c,h\rangle_n=0.
\]
Since $f_0=P_{\Vcal}y$, $\langle f_0-y,h\rangle_n=0$. Writing
$d_\rho=f_\rho-f_0\in\Vcal$ and $g=P_{\Vcal}c$ gives
\[
2d_\rho+\rho C_\rho g=0.
\]
Taking the inner product with $c$ yields
\[
C_\rho=C_0-\frac{\rho}{2}AC_\rho,
\]
which proves the result.
\end{proof}

\begin{proposition}[Accuracy cost]
\label{prop:path_cost}
Under Proposition~\ref{prop:path},
\[
\|f_\rho-y\|_n^2-\|f_0-y\|_n^2
=
\frac{C_0^2}{A}\bigl(1-q(\rho)\bigr)^2.
\]
\end{proposition}

Within a fixed prediction space, the targeted covariance falls to first order near $\rho=0$, while the in-sample squared-error cost is second order. For boosted trees this is a local benchmark only. Retraining can change splits, leaves, and stopping iteration.

\section{Linear and Boosting Implementations}
\label{app:implementation}


{
\color{red}

\subsection{Model-Specific Implementations}
\todo{This is from the main body}
\label{subsec:model_specific_implementations}

We now instantiate the two penalties for the model classes studied in Section~\ref{subsec:baseline_models}. Linear regression represents CCAO's earlier modeling framework and provides a transparent analytical benchmark. Second-order tree boosting represents the current LightGBM framework and the principal nonlinear implementation. These derivations also clarify the relevant computational distinction: the surrogate penalty is directly compatible with observation-wise custom-objective interfaces, whereas the direct penalty requires either a diagonal curvature approximation or a custom leaf-wise solver that retains its rank-one cross-observation curvature.

\subsubsection{Linear Covariance-Penalized Models}
\label{subsubsec:linear_covariance_models}

Consider the linear hypothesis class
\[
\mathcal H_{\mathrm{lin}}
=
\left\{
f_\theta:x\mapsto x^\top\theta
\;:\;
\theta\in\Theta\subseteq\mathbb R^p
\right\},
\]
where the feature vector may include an intercept. Let \(X\in\mathbb R^{n\times p}\) be the corresponding design matrix and write
\[
e=X\theta-y,
\qquad
c=Hy=y-\bar y\mathbf 1,
\qquad
H=I-\frac{1}{n}\mathbf 1\mathbf 1^\top,
\qquad
D=\diag(c_1^2,\ldots,c_n^2).
\]
Then \(C(f_\theta)=n^{-1}e^\top c\), and the baseline loss in~\eqref{eq:baseline_learning} is \(L(f_\theta)=n^{-1}\lVert X\theta-y\rVert_2^2\).

\paragraph{Direct squared-covariance penalty.}
For the direct formulation, the linear estimator solves
\begin{equation}
\label{eq:linear_direct_problem}
\theta_\rho^{\mathrm{cov}}
\in
\argminop_{\theta\in\Theta}
\left\{
\frac{1}{n}\lVert X\theta-y\rVert_2^2
+
\frac{\rho}{2n^2}
\bigl((X\theta-y)^\top c\bigr)^2
\right\}.
\end{equation}
Define
\[
a=X^\top c,
\qquad
b=y^\top c=c^\top c.
\]
When \(\Theta=\mathbb R^p\), the first-order optimality condition is
\begin{equation}
\label{eq:linear_direct_normal_equations}
\left(
X^\top X+\frac{\rho}{2n}aa^\top
\right)\theta
=
X^\top y+\frac{\rho}{2n}ba.
\end{equation}
Consequently, whenever the displayed inverse exists,
\begin{equation}
\label{eq:linear_direct_closed_form}
\theta_\rho^{\mathrm{cov}}
=
\left(
X^\top X+\frac{\rho}{2n}aa^\top
\right)^{-1}
\left(
X^\top y+\frac{\rho}{2n}ba
\right).
\end{equation}
Thus, the direct penalty produces a rank-one modification of the ordinary least-squares normal equations.

\paragraph{Sample-additive surrogate.}
For the surrogate penalty, the linear estimator solves
\begin{equation}
\label{eq:linear_surrogate_problem}
\theta_\rho^{\mathrm{surr}}
\in
\argminop_{\theta\in\Theta}
\left\{
\frac{1}{n}\lVert X\theta-y\rVert_2^2
+
\frac{\rho}{n}(X\theta-y)^\top D(X\theta-y)
\right\}.
\end{equation}
When \(\Theta=\mathbb R^p\), its first-order optimality condition is
\begin{equation}
\label{eq:linear_surrogate_normal_equations}
\left(X^\top X+\rho X^\top D X\right)\theta
=
X^\top y+\rho X^\top D y,
\end{equation}
and hence, whenever the displayed inverse exists,
\begin{equation}
\label{eq:linear_surrogate_closed_form}
\theta_\rho^{\mathrm{surr}}
=
\left(X^\top X+\rho X^\top D X\right)^{-1}
\left(X^\top y+\rho X^\top D y\right).
\end{equation}
Equivalently, the surrogate is weighted least squares with diagonal weight matrix \(I+\rho D\).

% \paragraph{Regularized linear models.}
% Both constructions extend directly to convex regularized estimators of the form
% \[
% \min_{\theta\in\Theta}
% \left\{
% L(f_\theta)+\Omega(\theta)+\rho\Psi(f_\theta)
% \right\}.
% \]
% A quadratic regularizer preserves a closed-form linear system. For a general differentiable convex regularizer, its gradient is added to the preceding optimality conditions. Nondifferentiable regularizers such as the LASSO preserve convexity but require an appropriate numerical method rather than a closed form.

\subsubsection{Covariance-Penalized Boosted Tree}
\label{subsubsec:covariance_tree_boosting}

We next specialize the penalties to the second-order boosting step in~\eqref{eq:baseline_boosting_step}. At iteration \(m\), let
\[
f_i=f^{(m-1)}(x_i),
\qquad
e_i=f_i-y_i,
\qquad
z=C\bigl(f^{(m-1)}\bigr)
=
\frac{1}{n}\sum_{i=1}^n e_i c_i.
\]
As in the standard boosting construction, the tree and its leaf weights are fitted to the current derivatives, after which the learning rate is applied in the ensemble update.

\paragraph{Direct penalty with diagonal curvature.}
Under the direct squared-covariance penalty, the gradient supplied to the boosting step is
\begin{equation}
\label{eq:boosting_direct_gradient}
g_i^{\mathrm{cov}}
=
\frac{2}{n}e_i
+
\frac{\rho}{n}z c_i.
\end{equation}
The exact Hessian is the dense matrix in~\eqref{eq:direct_gradient_hessian}, therefore it is not directly applicable in standard boosted tree packages. 

One heursitic is to retain only its diagonal, which gives
\begin{equation}
\label{eq:boosting_direct_diagonal_hessian}
h_i^{\mathrm{cov,diag}}
=
\frac{2}{n}
+
\frac{\rho}{n^2}c_i^2.
\end{equation}
Equations~\eqref{eq:boosting_direct_gradient}--\eqref{eq:boosting_direct_diagonal_hessian} define an observation-wise custom objective that can be supplied to standard second-order boosting software. The approximation preserves the exact gradient but ignores the cross-observation curvature generated by the direct penalty.



\paragraph{Sample-additive surrogate.}
For the surrogate objective, the observation-wise derivatives are
\begin{align}
\label{eq:boosting_surrogate_gradient_hessian}
g_i^{\mathrm{surr}}
&=
\frac{2}{n}e_i\left(1+\rho c_i^2\right),
\\
h_i^{\mathrm{surr}}
&=
\frac{2}{n}\left(1+\rho c_i^2\right).
\end{align}
Thus, once \(c_i^2\) has been precomputed, the surrogate can be implemented directly through the standard gradient and Hessian interface used by LightGBM. This observation-wise implementation is the formulation used for the principal practical specification in the CCAO application.

% \subsection{Interpretation and Scope of the Correction}
% \label{subsec:controls}

% The correction is deliberately targeted. The direct penalty controls the global linear relationship between log assessment ratios and log sale prices. The surrogate targets this relationship indirectly by increasing the loss assigned to residuals far from the center of the log-price distribution. Neither formulation is intended to replace the full set of assessor-facing diagnostics.

% The distinction between the training target and the evaluation measures is important. PRD, PRB, and VEI are computed on the price scale and have different functional forms, whereas the correction is constructed on the log scale. Appendix~\ref{app:cov_prd} shows that price-scale covariance is directionally aligned with PRD, but the training covariance is
% \[
% \Cov(\log r,\log P),
% \]
% not \(\Cov(r,P)\). A reduction in the training covariance is therefore expected to improve the dominant global ratio--value trend, but it does not algebraically guarantee a corresponding change in every vertical-equity statistic.

% The direct covariance penalty is also invariant to an additive shift in all log predictions. If \(e_i(f)\) is replaced by \(e_i(f)+a\) for a constant \(a\), then \(C(f)\) is unchanged. Equivalently, multiplying every price prediction by the same constant does not change the targeted covariance. The direct penalty therefore controls the slope of the ratio--value relationship but not the overall assessment level. The surrogate is not invariant to this shift because it penalizes \(e_i(f)^2\) rather than centered residuals. It may consequently affect assessment level, particularly in the tails, but it still does not impose exact aggregate calibration.

% More generally, the proposed correction does not guarantee:

% \begin{itemize}[leftmargin=6mm]
%     \item that PRD, PRB, or VEI will improve monotonically for every value of \(\rho\);
%     \item that the median, mean, or weighted mean assessment ratio will equal one;
%     \item that COD, COV, or predictive error will improve;
%     \item that nonlinear ratio--value patterns will be eliminated;
%     \item that countywide improvement will hold within every geographic area, property class, or time period; or
%     \item that a relationship estimated on the training sample will persist out of sample.
% \end{itemize}

% These limitations determine the evaluation protocol. The training covariance is treated as a mechanism measure, not as a sufficient model-selection criterion. The penalty level is selected on out-of-time validation data using the assessor-facing measures introduced in Section~\ref{subsec:metrics}. Predictive accuracy, assessment level, horizontal uniformity, vertical equity, and subgroup behavior are then evaluated separately. In this way, the correction moves a specific vertical-equity concern into model training while retaining the broader ratio-study process needed to assess its consequences.





}

\subsection{Linear Model}

Let $X\in\R^{n\times p}$ and $e=X\theta-y$. With
\[
c=y-\bar y\1,\qquad a=X^\top c,\qquad b=y^\top c,
\]
the direct penalized estimator under loss $\|X\theta-y\|_2^2/n$ satisfies
\[
\left(
\frac{1}{n}X^\top X+\frac{\rho}{2n^2}aa^\top
\right)\theta
=
\frac{1}{n}X^\top y+\frac{\rho}{2n^2}ba.
\]
This is a rank-one modification of the normal equations.

For the surrogate, let $D=\diag(c_1^2,\ldots,c_n^2)$. Then
\[
\left(
\frac{1}{n}X^\top X+\frac{\rho}{n}X^\top DX
\right)\theta
=
\frac{1}{n}X^\top y+\frac{\rho}{n}X^\top Dy.
\]

\subsection{Direct Penalty in Boosting}

For current predictions $f$, define $e=f-y$ and
\[
z=\frac{1}{n}e^\top c.
\]
Under squared loss, the direct objective has
\[
g_i=\frac{2e_i}{n}+\frac{\rho}{n}zc_i,
\]
and exact Hessian
\[
\nabla_f^2
=
\frac{2}{n}I+\frac{\rho}{n^2}cc^\top.
\]
The diagonal approximation uses
\[
h_i^{\mathrm{diag}}
=
\frac{2}{n}+\frac{\rho}{n^2}c_i^2.
\]
For a fixed tree with leaves $I_1,\ldots,I_J$, let
\[
b_j=\sum_{i\in I_j}e_i,
\qquad
s_j=\sum_{i\in I_j}c_i,
\qquad
n_j=|I_j|.
\]
With leaf-weight vector $w$, ridge parameter $\lambda$, and $s=(s_1,\ldots,s_J)^\top$, the exact local quadratic system is
\[
A_Tw^\star=-g_T,
\qquad
g_{T,j}=\frac{2b_j}{n}+\frac{\rho z s_j}{n},
\]
\[
A_T
=
\diag\!\left(\frac{2n_1}{n}+\lambda,\ldots,\frac{2n_J}{n}+\lambda\right)
+\frac{\rho}{n^2}ss^\top.
\]
Thus $w^\star=-A_T^{-1}g_T$ when $A_T$ is nonsingular. This solver retains the cross-observation curvature but requires a custom tree-construction procedure; it is not available through the usual observation-wise API.

\subsection{Surrogate in Boosting}

For
\[
\Psi^{\mathrm{surr}}(f)
=
\frac{1}{n}\sum_i e_i^2c_i^2,
\]
the derivatives are observation-wise:
\[
g_i=\frac{2e_i}{n}+\frac{2\rho}{n}e_ic_i^2,
\qquad
h_i=\frac{2}{n}+\frac{2\rho}{n}c_i^2.
\]
This is the implementation used for the primary practical model.

\section{CCAO Reproducibility Requirements}
\label{app:ccao_results}

\advisorchange{A replication package for the CCAO application should contain} the penalty grid, fold-level validation predictions, active selection safeguards, selected configuration, final test predictions, and assessor metrics before and after any separate level calibration. It should include both the ordinary LightGBM baseline and a custom-objective fit with zero penalty. Before operational use, the same analysis should report block-bootstrap intervals and ratio curves by township, property class, value range, and time. These outputs separate model-selection uncertainty, implementation effects, and local variation.

\section{Exploratory ATTOM Benchmark and Recalibration}
\label{app:attom}

The exploratory ATTOM benchmark uses single-family recorder transfers \advisorchange{that pass the benchmark price and data-quality filters}. Each transfer is joined to the most recent assessor-history record available before the sale and to lagged contextual variables. The final 20\% of eligible sales form a chronological test block, with a validation block inside the remaining data.

The benchmark compares the ordinary LightGBM fit, the custom-objective zero-penalty control, positive covariance-penalty values, and centered recalibration maps of the baseline log predictions. Recalibration parameters are estimated without using the test block. Constant level calibration is reported separately because it changes the overall ratio level but not the targeted log-residual covariance.

The current county analysis was developed iteratively and is treated as exploratory. A confirmatory version must freeze the feature set, LightGBM settings, penalty grid, recalibration family, and selection rule before opening any county's test block. It must report the full non-dominated frontier for prediction error, PRD, PRB, COD, assessment level, and value-decile ratio curves. Until that protocol is rerun, the county results are evidence about possible behavior and failure modes rather than definitive external validation.

\section{Alternative Selection and Stacking Extensions}
\label{app:selection_extensions}

When no candidate satisfies every active validation safeguard, the office should first inspect the source of infeasibility rather than silently replacing the constraints with a weighted score. Possible actions are to improve the data, revise the candidate set, separate level calibration from vertical-trend control, or report the least-violating candidate without calling it feasible.

A normalized distance-to-ideal rule can be used for exploratory comparison. Convex stacking can also enlarge the attainable set when constraints such as PRD can be written as affine inequalities in the stacking weights. These methods are secondary because they add choices and can make the final model harder to explain. The discrete constrained rule in \eqref{eq:selection} remains the recommended default.

\section*{Data, Code, and Institutional Statement}

\advisorchange{The primary application uses CCAO data} and a translated version of the office's published residential modeling workflow. The multi-county benchmark uses licensed ATTOM-derived records that cannot be redistributed through the paper. Reproducibility therefore depends on archiving the custom objectives, preprocessing specifications, temporal splits, and evaluation code, together with instructions for authorized data users. The results reported here are research findings, not official CCAO assessment decisions. The paper makes no claim that CCAO has adopted the correction or placed it into production.




\end{document}
